\documentclass[12pt]{article}
%\usepackage{amssymb,epsf,graphicx}
\usepackage{amsmath,amssymb,amsthm, bm,epsf,epsfig,graphicx,hyperref,physics, cleveref, subfig, slashed, stmaryrd, tensor, mathrsfs,}

\usepackage[backend=bibtex,natbib,style=numeric-comp,sorting=none, doi=false, isbn=false,url=false]{biblatex}
\usepackage{array,booktabs}
\input epsf.sty
\topmargin -.5cm \textheight 21cm \oddsidemargin -.125cm
\textwidth 17cm
\setlength{\parskip}{0.5em}

\usepackage[]{todonotes}

\usepackage[weather, alpine]{ifsym}

\definecolor{bisque}{rgb}{1.0, 0.89, 0.77}
\definecolor{forestgreen(web)}{rgb}{0.13, 0.55, 0.13}

\newcommand{\ie}{\begin{equation}\begin{aligned}}
\newcommand{\fe}{\end{aligned}\end{equation}}
\numberwithin{equation}{section}
\newcommand{\R}{\mathrm{R}}
\newcommand{\NS}{\mathrm{NS}}
\newcommand{\bpzbraket}[2]{%
  \bigl\langle\!\bigl\langle #1 \bigm| #2 \bigr\rangle%
}

\renewcommand{\title}[1]{\vbox{\center\LARGE{#1}}\vspace{5mm}}
\renewcommand{\author}[1]{\vbox{\center#1}\vspace{5mm}}
\newcommand{\address}[1]{\vbox{\center\em#1}}
\newcommand{\email}[1]{\vbox{\center\tt#1}\vspace{5mm}}
\newcommand{\WT}[1]{\widetilde{#1}}

\newtheorem{proposition}{Proposition}

\begin{document}

This is a working notes focusing on deriving the $c$-recursion relation for 2D $\mathcal{N}=1$ superconformal blocks in arbitrary plumbing channel. Written by AI (an individual), not by AI (artificial intelligence).



\section{Virasoro conformal blocks}

We first review the recipe for Virasoro conformal blocks, following \cite{Cho-Collier-Yin}, with the genus-2 conformal block in the theta channel as an example. The derivation can be generalized to arbitrary genus and arbitrary plumbing channel. In this case, the Virasoro conformal block is defined as
\ie
F\left(h_1, h_2, h_3, c\right)= & \sum_{|A|=|B|,|C|=|D|,|E|=|F|} q_1^{|A|} q_2^{|C|} q_3^{|E|} G_{h_1}^{A B} G_{h_2}^{C D} G_{h_3}^{E F} \\
& \times \rho\left(L_{-A} \nu_1, L_{-C} \nu_2, L_{-E} \nu_3\right) \rho\left(L_{-B} \nu_1, L_{-D} \nu_2, L_{-F} \nu_3\right) ,
\fe
where $A$ is a partition of $n=|A|$ in descending order, $L_{-A}$ stands for a string of lowering operators made of $L_{-m}$'s, $G^{AB}_h$ is the inverse of the Gram matrix $\langle h|L_{-A}^\dagger L_{-B} |h \rangle$, and $\rho$ is the 3 point function of Virasoro descendents of the canonically normalized primaries $\nu_i$ with the structure constant stripped off:
\ie \label{rho}
\langle\langle\xi_3 \otimes \bar{\xi}_3| \xi_2\otimes \bar{\xi}_2(z, \bar{z})|\xi_1 \otimes \bar{\xi}_1\rangle=C_{321} \rho\left(\xi_3, \xi_2, \xi_1| z\right) \rho\left(\bar{\xi}_3, \bar{\xi}_2, \bar{\xi}_1 | \bar{z}\right).
\fe
We will also abbreviate $\rho(\xi_1, \xi_2, \xi_3|1)$ as $\rho(\xi_1, \xi_2, \xi_3)$.

The stretegy of computing this conformal block is to treat it as a meromorphic function in $c$ at fixed $h_i$. The poles on the complex $c$-plane can be located at the values where the Gram matrix develops a null state, and the whole conformal block can be determined by its large-$c$ behavior and the residue at these poles. Namely,
\ie
F(h_1,h_2, h_3,c) = F(h_1,h_2, h_3,c\rightarrow \infty) + \sum_{i, r\geq 2, s\geq 1} \frac{W_{rs}^{(i)}(h_1, h_2, h_3, c)}{c-c_{rs}(h_i)}.
\fe
Note that we are considering the case where all the internal weights $h_i$ are different, hence there are no double poles from two internal multiplets admits a null state simutaneously.

The poles at $c=c_{rs}(h_i)$ arises when the $i$-th intermediate multiplet develops a null state. The Kac formula then tells us
\ie
c_{r s}(h)=1+6\left(b_{r s}(h)+b_{r s}(h)^{-1}\right)^2,\quad  b_{r s}(h)^2=\frac{r s-1+2 h+\sqrt{ (r-s)^2+4(r s-1) h+4 h^2}}{1-r^2}
\fe
where $r=2,3,4, \cdots, s=1,2,3, \cdots$. Conversely, let's denote the weight $h$ that develops a null state at level $rs$ at fixed $c$ by $d_{rs}(c)$. At this value of $c$, the multiplet with lowest weight state $\nu_{d_{rs}}$ admits a null state $\chi_{rs}=\sum_{|M|=rs}\chi_{rs}^M L_{-M}\nu_{d_{rs}}$. In this case, the residue $W_{rs}^{(i)}$ receives only contributions from the descendants of this null state. For example when $i=1$, one have
\ie
W^{(1)}_{rs} &= -\frac{\partial c_{rs}(h_1)}{\partial h_1}\lim_{h_1\rightarrow d_{r,s}} \frac{h_1-d_{rs}}{\|\sum_{|M|=rs} \chi_{rs}^M L_{-M}\nu_h \|^2} \\
&\times  \sum_{|A|=|B|,|C|=|D|,|E|=|F|} q_1^{|A|+rs} q_2^{|C|} q_3^{|E|} G_{h_1+rs}^{A B} G_{h_2}^{C D} G_{h_3}^{E F} \\
&\quad \quad \quad  \times \rho\left(L_{-A} \chi_{rs}, L_{-C} \nu_2, L_{-E} \nu_3\right) \rho\left(L_{-B}  \chi_{rs}, L_{-D}\nu_2, L_{-F} \nu_3\right).
\fe
Using the factorization relation
\ie
\rho\left(L_{-N} \chi_{r s}, L_{-M} \nu_2, L_{-P} \nu_3 \right) & =\rho\left(L_{-N} \nu_{d_{r s}+r s}, L_{-M} \nu_2, L_{-P} \nu_3 \right) \rho\left(\chi_{r s}, \nu_2, \nu_3 \right) \\
\rho\left(L_{-N} \nu_1, L_{-M} \chi_{r s}, L_{-P} \nu_3 \right) & =\rho\left(L_{-N} \nu_1, L_{-M} \nu_{d_{r s}+r s}, L_{-P} \nu_3 \right) \rho\left(\nu_1, \chi_{r s}, \nu_3 \right) \\
\rho\left(L_{-N} \nu_1, L_{-M} \nu_2, L_{-P} \chi_{r s} \right) & =\rho\left(L_{-N} \nu_1, L_{-M} \nu_2, L_{-P} \nu_{d_{r s}+r s} \right) \rho\left(\nu_1, \nu_2, \chi_{r s} \right) \\
\rho\left(L_{-N} \chi_{r s}, L_{-M} \nu_2, L_{-P} \chi_{r s} \right) & =\rho\left(L_{-N} \nu_{d_{r s}+r s}, L_{-M} \nu_2, L_{-P} \nu_{d_{r s}+r s} \right) \rho\left(\chi_{r s}, \nu_2, \chi_{r s} \right) ,
\fe
one can then arrive at
\ie
W^{(1)}_{rs} &= -\frac{\partial c_{rs}(h_1)}{\partial h_1} A^c_{rs} q_1^{rs} [P_c^{rs}(h_3, h_2)]^2 F(h_1\rightarrow h_1+rs,c\rightarrow c_{rs}(h_1)),
\fe
where the quantity $A^c_{rs}$ and the fusion polynomial $P_c^{rs}(h_3, h_2)$ are defined as
\ie
&A_{r s}^c=\lim _{h \rightarrow d_{r s}}\frac{h-d_{r s}}{\|\sum_{|M|=rs} \chi_{rs}^M L_{-M}\nu_h \|^2}=\frac{1}{2} \prod_{m=1-r}^r \prod_{n=1-s}^s\left(m b_{rs}+n b^{-1}_{rs}\right)^{-1}, \quad(m, n) \neq(0,0),(r, s)\\
&P_c^{r s}(h_1, h_2)  \equiv \rho\left(\chi_{r s}, \nu_2, \nu_1\right) =\prod_{p=1-r \text { step } 2}^{r-1} \prod_{q=1-s \text { step } 2}^{s-1} \frac{\lambda_1+\lambda_2+p b_{rs}+q b^{-1}_{rs}}{2} \frac{\lambda_1-\lambda_2+p b_{rs}+q b^{-1}_{rs}}{2},
\fe
where the products for the fusion polynomial are taken over $p+r=1 (\text{mod }2)$ and $q+s=1 (\text{mod }2)$, and $\lambda_i$ are defined by $h_i = \frac{1}{4}(b_{rs}+b^{-1}_{rs})^2-\frac{1}{4}\lambda_i^2$. The other residues $W_{rs}^{(2)}, W_{rs}^{(3)}$ follows similarly. 

The large-$c$ limit $F(h_1,h_2, h_3,c\rightarrow \infty)$ of the conformal block can be computed by the following. Denoting by $\langle A\rangle$ the number of non-$L_{-1}$ generators in $L_{-A}$. The key observation is that, the Gram matrix element $G_{AB}$ scales like $c^{\langle A\rangle}$ for $A=B$, no faster than $c^{\langle A\rangle-1}$ for $\langle A\rangle = \langle B\rangle$ but $A\neq B$, and no faster than $c^{\text{min}(\langle A\rangle, \langle B\rangle)}$ for $\langle A\rangle \neq \langle B\rangle$. This can be directly seen from the Virasoro algebra, where the leading term in $c$ is given by the commutators $[L_m,L_{-m}]$. One can then start from the basis of states $\{L_{-A}\nu, |A|=n\}$ at level $n$, apply Gram-Schmidt orthogonalization, and obtain a basis of orthogonal states $\mathcal{L}_{-A}\nu$. These states are of the form
\ie 
\mathcal{L}_{-A}\nu = L_{-A}\nu + \sum_{|B|=n, \langle B\rangle \leq \langle A\rangle, B\neq A} g_{AB}(c,h) L_{-B}\nu,
\fe
where $g_{AB}(c,h)$ is order $\mathcal{O}(c^0)$ for $B < A$ and $\mathcal{O}(c^{-1})$ for $\langle B\rangle = \langle A\rangle$. In this basis, the conformal block can be written as
\ie
F(h_1, h_2, h_3, c) = \sum_{A, B, C} q_1^{|A|}q_2^{|B|}q_3^{|C|} \frac{\rho(\mathcal{L}_{-A}\nu_1, \mathcal{L}_{-B}\nu_2, \mathcal{L}_{-C}\nu_3)^2}{\|\mathcal{L}_{-A}\nu_1\|^2\|\mathcal{L}_{-B}\nu_2\|^2\|\mathcal{L}_{-C}\nu_3\|^2}
\fe
where $\|\mathcal{L}_{-A}\nu_1\|^2$ is of order $c^{\langle A\rangle}$.

Now let's consider the scaling behavior of $\rho(\mathcal{L}_{-A}\nu_1, \mathcal{L}_{-B}\nu_2, \mathcal{L}_{-C}\nu_3)$ in the $c\rightarrow \infty$ limit. In this limit, only the term of order $c^{(\langle A\rangle+\langle B\rangle+\langle C\rangle)/2}$ survives in the conformal block. For a general 3-point function $\rho(L_{-A}\nu_1, L_{-B}\nu_2, L_{-C}\nu_3)$, the leading term must come from pairwise contractions $[L_m, L_{-m}]$ of non-$L_{-1}$ Virasoro generators, and the maximum number of such contractions is precisely $(\langle A\rangle+\langle B\rangle+\langle C\rangle)/2$, which gives a scaling behavior of $c^{(\langle A\rangle+\langle B\rangle+\langle C\rangle)/2}$. In each term of $\mathcal{L}_{-A}\nu$, only the $L_{-A}\nu$ term can give this scaling, since $g_{AB}(c,h)$ is subleading for $\langle B\rangle = \langle A\rangle$ and the correlator is subleading for $\langle B\rangle<\langle A\rangle $. Moreover, since only the central term in the commutator contributes to this limit, we have
\ie
\rho(\mathcal{L}_{-A}\nu_1, \mathcal{L}_{-B}\nu_2, \mathcal{L}_{-C}\nu_3)\rightarrow \rho(L_{-A'}\mathbf{1}, L_{-B'}\mathbf{1}, L_{-C'}\mathbf{1})\rho(L_{-1}^{k_A}\nu_1, L_{-1}^{k_B}\nu_2, L_{-1}^{k_C}\nu_3),
\fe
where $\mathbf{1}$ is the identity operator, and we have written $A= A'\cup \{1\}^{k_A}$. Similarly, the 2-point function also admits a similar factorization relation at large $c$, namely
\ie
\|\mathcal{L}_{-A}\nu\|^2 \rightarrow \|L_{-A'}\mathbf{1}\|^2 \|L_{-1}^{k_A}\nu\|^2.
\fe
Plugging these relations back to the conformal block, one can then arrive at the large-$c$ limit
\ie
F(h_1, h_2, h_3,c\rightarrow \infty) = F(0, 0, 0,c\rightarrow \infty)\times F_{\mathfrak{sl}(2)}(h_1, h_2, h_3),
\fe
where the first factor is the large-$c$ limit of the vacuum block, and the second factor is the global $\mathfrak{sl}(2)$ block that's independent of $c$.  The global $\mathfrak{sl}(2)$ blocks are generally simple and each term in the $q$-sum can be computed in terms of Pochhammer symbols (See (4.19)-(4.22) of \cite{Cho-Collier-Yin}), while the vacuum block in the Schottky frame has a simple product formula (See (5.5)).

Assembling all the pieces, we arrive at the final form for the $c$-recursion relation for genus-2 Virasoro conformal blocks in the theta channel:
\ie
F(h_1, h_2, h_3, c) &= F(0, 0, 0,c\rightarrow \infty)\times F_{\mathfrak{sl}(2)}(h_1, h_2, h_3)\\
&+ \sum_{r\geq 2,s\geq 1} -\left(\frac{\partial c_{rs}(h_1)}{\partial h_1}\right) \frac{A^c_{rs} q_1^{rs} [P_c^{rs}(h_3, h_2)]^2}{c-c_{r,s}(h_1)} F(h_1\rightarrow h_1+rs,c\rightarrow c_{rs}(h_1)) \\
&+ 2\text{ permutation in }h_1, h_2, h_3.
\fe
Using this relation, one can recursively evalute the conformal block $F(h_1, h_2, h_3, c)$ as a $q$-expansion, as the polar terms are all higher-order in $q$ due to the factor $q^{rs}$.

\section{$\mathcal{N}=1$ superconformal algebra}

Let \(\nu\in\left\{0,\frac12\right\}\). The two-dimensional \(\mathcal N=1\) super-Virasoro algebra \(\mathsf{SCA}_\nu\) is generated by
\[
\{L_n,G_r\mid n\in\mathbb Z,\ r\in\mathbb Z+\nu\}.
\]
Its nonvanishing (anti)commutation relations are
\ie
\left[L_m,L_n\right]
&=(m-n)L_{m+n}
+\frac{c}{12}\left(m^3-m\right)\delta_{m+n,0},\\
\left[L_m,G_r\right]
&=\left(\frac{m}{2}-r\right)G_{m+r},\\
\left\{G_r,G_s\right\}
&=2L_{r+s}
+\frac{c}{3}\left(r^2-\frac14\right)\delta_{r+s,0}.
\fe
Here, \(\nu=0\) and \(\nu=\frac12\) correspond to the Ramond and Neveu–Schwarz sectors, respectively.

A two-dimensional \(\mathcal N=(1,1)\) superconformal field theory contains both holomorphic and anti-holomorphic superconformal algebras. We denote the anti-holomorphic algebra by \(\widetilde{\mathsf{SCA}}_{\tilde\nu}\), with generators
\[
\{\widetilde L_n,\widetilde G_r\mid n\in\mathbb Z,\ r\in\mathbb Z+\tilde\nu\}.
\]
They obey
\ie
\left[\widetilde L_m,\widetilde L_n\right]
&=(m-n)\widetilde L_{m+n}
+\frac{\widetilde c}{12}\left(m^3-m\right)\delta_{m+n,0},\\
\left[\widetilde L_m,\widetilde G_r\right]
&=\left(\frac{m}{2}-r\right)\widetilde G_{m+r},\\
\left\{\widetilde G_r,\widetilde G_s\right\}
&=2\widetilde L_{r+s}
+\frac{\widetilde c}{3}\left(r^2-\frac14\right)\delta_{r+s,0}.
\fe
The holomorphic and anti-holomorphic generators mutually (anti)commute.


\section{Superconformal Block Decomposition}

We now generalize this $c$-recursion relation to superconformal blocks in $\mathcal{N}=1$ superconformal field theories. The sphere 4-point case of the recursion relation has been worked out in \cite{0611266, 0611295}, and the generalization to general plumbing channel has been worked out in \cite{1806.09563}.

Apart from the bosonic moduli, one also needs to specify a spin structure to put a SCFT on a Riemann surface. A plumbing channel is specified by a choice of NS and R for each punctures, a modulus $q$ and a sign $\eta$ for each pair of plumbed punctures. Only NS and NS, R and R punctures are allowed to be plumbed together, and the NS/R configuration of punctured on each 3-punctured sphere must be (NS, NS, NS) or (NS, R, R). For genus-$g$ surface, there are $2^g$ choices of NS/R punctures and $2^g$ independent sign choices for each tube.


Similar to the notation $L_{-A}\nu$, we will denote a NS superconformal primary by $\phi$, and its descendent by
\ie
\mathbf{L}_{-A}\phi \equiv L_{-n_1}\cdots L_{-n_\ell} G_{-r_1}\cdots G_{-r_p} \phi,\quad |A|\equiv \sum_{i=1}^\ell n_i + \sum_{j=1}^{p} r_j,\quad (-1)^A \equiv (-1)^{p}
\fe
where $n_1 \geq ... \geq n_\ell, r_1>...>r_p$. We will also take $\phi$ to be Grassmann even. \todo[inline]{Yutai: I do not know why we can choose the parity of phi in general. I would assume it to be generic.} 

In the language of 2D $\mathcal{N}=1$ superspace, one can pack $\phi$ and its level-$\frac{1}{2}$ descendent $\chi \equiv G_{-1/2}\phi$ in the primary superfield $\Phi(z,\theta)$, defined by
\ie
\Phi(z,\theta) = \phi(z)+\theta \chi(z).
\fe
Superconformal symmetry then constraints the 3-point functions to be of the following form
\ie \label{superspace23pt}
& \left\langle\Phi_1(z_1, \theta_1) \Phi_2(z_2, \theta_2)\right\rangle \propto \frac{\delta_{h_1 h_2}}{Z_{12}^{2 h_1}}, \\
& \left\langle\Phi_1(z_1, \theta_1) \Phi_2(z_2, \theta_2) \Phi_3(z_3, \theta_3)\right\rangle=\frac{\left\langle\Phi_1(0,0) \Phi_2(1, \Theta) \Phi_3^{\prime}(\infty, 0)\right\rangle}{Z_{21}^{h_1+h_2-h_3} Z_{31}^{h_1+h_3-h_2} Z_{32}^{h_2+h_3-h_1}},
\fe
where
\ie
& Z_{i j} \equiv z_i-z_j-\theta_i \theta_j \\
& \Theta \equiv \frac{\theta_1 z_{23}+\theta_2 z_{31}+\theta_3 z_{12}-\frac{1}{2} \theta_1 \theta_2 \theta_3}{\sqrt{z_{12} z_{13} z_{23}}} .
\fe
Importantly, there are now two independent structure constants for each triple of NS supermultiplets, they are the $1$ and $\Theta$ component of the 3-point function, given essentially by $C_{123}=\langle \phi_1(0)\phi_2(1)\phi_3(\infty)\rangle$ and $\tilde{C}_{123}=\langle \phi_1(0)\chi_2(1)\phi_3(\infty)\rangle$. This is due to the fact that the 3-NS punctured genus-0 super-Riemann surface has 1 odd moduli, which is given precisely by this odd cross-ratio $\Theta$.

For 3-point functions of superconformal descendants $\langle\mathbf{L}_{-A}\phi_1\mathbf{L}_{-B}\phi_2\mathbf{L}_{-C}\phi_3\rangle$, It can be shown by superconformal Ward identity that the $|A|+|B|+|C|\in \mathbb{Z}$ case only depends on $C_{123}$, while the $|A|+|B|+|C|\in \mathbb{Z}+\frac12$ case only depends on $\tilde{C}_{123}$. We can then define two types of normalized 3-point functions
\ie 
\langle\langle\xi_3 \otimes \bar{\xi}_3| \xi_2\otimes \bar{\xi}_2(z, \bar{z})|\xi_1 \otimes \bar{\xi}_1\rangle = (\pm)\begin{cases}
C_{321} \rho_0(\xi_3, \xi_2, \xi_1| z) \rho_0(\bar{\xi}_3, \bar{\xi}_2, \bar{\xi}_1 | \bar{z}), \ \  \text{integer total level},\\
\tilde{C}_{321} \rho_1(\xi_3, \xi_2, \xi_1| z) \rho_1(\bar{\xi}_3, \bar{\xi}_2, \bar{\xi}_1 | \bar{z}),\ \  \text{half-integer total level},
\end{cases}
\fe
by essentially normalizing 
\ie
\rho_0(\phi_1, \phi_2, \phi_3|1)=1,\quad \rho_1(\phi_1, \chi_2, \phi_3|1)=1.
\fe 
The sign $(\pm)$ is here indicating that we use the operator order $\langle \xi_3\xi_2\xi_1\bar{\xi}_3\bar{\xi}_2\bar{\xi}_1 \rangle$ to define the $\rho$ on the RHS, which can differ from the operator order in LHS by a sign. Again, we denote $\rho_a(\xi_1, \xi_2, \xi_3)=\rho_a(\xi_1, \xi_2, \xi_3|1)$.

The anti-holomorphic sector is entirely analogous to the holomorphic sector. We distinguish quantities in the anti-holomorphic sector by adding a tilde over them.
The standard BPZ bilinear pairing in a superconformal field theory would factorize into holomorphic and antiholomorphic pieces as
\ie
\bpzbraket{ \mathbf{L}_{-A} \phi  \otimes \widetilde{\mathbf{L}}_{-\widetilde{A}} \widetilde{\phi} }{ \mathbf{L}_{-B} \phi \otimes  \widetilde{\mathbf{L}}_{-\widetilde{B}} \widetilde{\phi}} =(-1)^{(B+p_{\phi})(\widetilde{A}+p_{\widetilde{\phi}})}\bpzbraket{\mathbf{L}_{-A} \phi}{\mathbf{L}_{-B} \phi } \times \bpzbraket{ \widetilde{\mathbf{L}}_{-\widetilde{A} }\widetilde{\phi}}{ \widetilde{\mathbf{L}}_{-\widetilde{B}}\widetilde{\phi}}.
\fe


\subsection{Conformal Block Decomposition in NS Plumbing}

For illustration, we consider a 2d \(\mathcal N=(1,1)\) superconformal field theory and derive the conformal-block decomposition of its genus-2 partition function. The final result present a nice holomorphic factorialization structure after keeping track of the signs in the computation.

Let $\mathcal{O}_i(z,\bar{z})$ be the superconformal primary operator with conformal weight $(h_i,\tilde{h}_i)$. In the following derivation, for convenience, we will assume $\mathcal{O}_i$ to be orthogonal to each other, with
\ie
 \bpzbraket{O_i }{O_j} =D_i \delta_{ij}.
\fe
Here $\left\langle\!\left\langle \psi \big| \right.\right.$ denotes the BPZ conjugate of state $\psi$. 

Primary operator $\mathcal{O}_i(z,\bar{z})$ admits a formal factorialization 
\ie
\mathcal{O}_i(z,\bar{z}) = \sqrt{(-1)^{p_i \tilde{p}_i}D_i} \phi_i (z)\otimes \widetilde \phi_i(\bar{z}),
\fe
with $\phi$ the holomorphic superconformal primary with parity $p_i$ and  $\widetilde \phi(\bar{z})$ the anti-holomorphic superconformal primary with parity $\tilde{p}_i$, normalized such that 
\ie
\bpzbraket{\phi }{\phi} =\bpzbraket{\widetilde{\phi} }{\widetilde \phi} =1.
\fe
The parity of $\mathcal{O}_i$ is $p_i+\tilde{p}_i$. In the following, we will derive the conformal block decomposition for genus-two partition function \(Z_{g=2}\) in theta channel and glass channel. 

\subsubsection*{All-NS Block Decomposition in Theta Channel}

Consider the genus-two partition function \(Z_{g=2}\). We work in the theta-graph plumbing channel, taking all three tubes to have (NS,NS) spin structure in both the holomorphic and anti-holomorphic sectors. In the theta plumbing channel, the genus-two partition function \(Z_{g=2}\) has a prescription
\ie
Z^{\Theta}_{g=2}
={}&
\sum_{\mathcal{O}_1,\mathcal{O}_2,\mathcal{O}_3}
\sum_{\substack{
|A|=|B|,\ |C|=|D|,\ |E|=|F|\\
|\widetilde A|=|\widetilde B|,\ 
|\widetilde C|=|\widetilde D|,\ 
|\widetilde E|=|\widetilde F|
}}
\prod_{i=1}^{3}
q_i^{h_i}\widetilde q_i^{\widetilde h_i} \eta_i^{p_i} \widetilde \eta_i^{\tilde{p}_i}
\times
q_1^{|A|}\eta_1^A
q_2^{|C|}\eta_2^C
q_3^{|E|}\eta_3^E
\widetilde q_1^{|\widetilde A|}
\widetilde\eta_1^{\widetilde A}
\widetilde q_2^{|\widetilde C|}
\widetilde\eta_2^{\widetilde C}
\widetilde q_3^{|\widetilde E|}
\widetilde\eta_3^{\widetilde E}
\\
&\times
(-1)^{
(A+\widetilde A +p_1+\tilde{p}_1)(C+\widetilde C+p_2+\tilde{p}_2)
+(A+\widetilde A+p_1+\tilde{p}_1)(E+\widetilde E+p_3+\tilde{p}_3)
+(C+\widetilde C+p_2+\tilde{p}_2)(E+\widetilde E+p_3+\tilde{p}_3)}
\\
&\times
(\mathbb{G}_1)^{A\widetilde A;B\widetilde B}
(\mathbb{G}_2)^{C\widetilde C;D\widetilde D}
(\mathbb{G}_3)^{E\widetilde E;F\widetilde F}
\\
&\times
\left\langle\!\left\langle
\mathbf{L}_{-A} \widetilde{\mathbf{L}}_{-\widetilde A}\mathcal{O}_1
\middle|
\mathbf{L}_{-C} \widetilde{\mathbf{L}}_{-\widetilde C}\mathcal{O}_2(1,1)
\middle|
\mathbf{L}_{-E} \widetilde{\mathbf{L}}_{-\widetilde E}\mathcal{O}_3
\right\rangle\right.
\times
\left\langle\!\left\langle
\mathbf{L}_{-B} \widetilde{\mathbf{L}}_{-\widetilde B}\mathcal{O}_1
\middle|
\mathbf{L}_{-D} \widetilde{\mathbf{L}}_{-\widetilde D}\mathcal{O}_2(1,1)
\middle|
\mathbf{L}_{-F} \widetilde{\mathbf{L}}_{-\widetilde F}\mathcal{O}_3
\right\rangle\right.
\fe
In this prescription, we need to choose the Weyl frame to be the one defined by theta channel plumbing. The sign in the sum comes from reordering the operators into the prescription of the matrix element. The gram matrix is defined by
\ie
(\mathbb{G}_i)_{ A\widetilde A;B\widetilde B}
&\equiv
\bpzbraket{
\mathbf{L}_{-A}
\widetilde{\mathbf{L}}_{-\widetilde A}
\mathcal{O}_i
}{
\mathbf{L}_{-B}
\widetilde{\mathbf{L}}_{-\widetilde B}
\mathcal{O}_i
}
\fe
It can be factorized into the holomorphic and anti-holomorphic pieces as
\ie\label{eq: Gram_Factorialization}
(\mathbb{G}_i)_{ A\widetilde A;B\widetilde B}&= D_i
(-1)^{(\widetilde A+ \tilde{p}_i)  (B+p_i)+ p_i\tilde{p}_i}
(\mathbf{G}_{h_i})_{AB}
(\widetilde{\mathbf{G}}_{\widetilde h_i})_
{\widetilde A\widetilde B}.
\fe
where
\ie
(\mathbf{G}_{h_i})_{AB} : = \bpzbraket{\mathbf{L}_{-A}\phi_i }{\mathbf{L}_{-B}\phi_i},~ (\widetilde{\mathbf{G}}_{h_i})_{\widetilde{A}\widetilde{B}} : = \bpzbraket{\widetilde{\mathbf{L}}_{-\widetilde{A}}\tilde{\phi}_i }{\widetilde{\mathbf{L}}_{-\widetilde B}\tilde{\phi}_i}.
\fe
The three point function would also factorize
\ie
&\left\langle\!\left\langle
\mathbf{L}_{-A}
\widetilde{\mathbf{L}}_{-\widetilde A}\mathcal{O}_1
\middle|
\mathbf{L}_{-C}
\widetilde{\mathbf{L}}_{-\widetilde C}\mathcal{O}_2(1,1)
\middle|
\mathbf{L}_{-E}
\widetilde{\mathbf{L}}_{-\widetilde E}\mathcal{O}_3
\right\rangle\right.
\\
&\quad\times
\left\langle\!\left\langle
\mathbf{L}_{-B}
\widetilde{\mathbf{L}}_{-\widetilde B}\mathcal{O}_1
\middle|
\mathbf{L}_{-D}
\widetilde{\mathbf{L}}_{-\widetilde D}\mathcal{O}_2(1,1)
\middle|
\mathbf{L}_{-F}
\widetilde{\mathbf{L}}_{-\widetilde F}\mathcal{O}_3
\right\rangle\right.
\\
={}&
(-1)^{p_1\tilde{p}_1 +p_2\tilde{p}_2+ p_3\tilde{p}_3}
D_1D_2D_3
\left(C_{123}^{(a)}\right)^2
\\
&\quad\times
\rho_a\!\left(
\mathbf{L}_{-A}\phi_1,
\mathbf{L}_{-C}\phi_2,
\mathbf{L}_{-E}\phi_3
\right)
\widetilde{\rho}_{\tilde{a}}\!\left(
\widetilde{\mathbf{L}}_{-\widetilde A}\widetilde{\phi}_1,
\widetilde{\mathbf{L}}_{-\widetilde C}\widetilde{\phi}_2,
\widetilde{\mathbf{L}}_{-\widetilde E}\widetilde{\phi}_3
\right)
\\
&\quad\times
\rho_a\!\left(
\mathbf{L}_{-B}\phi_1,
\mathbf{L}_{-D}\phi_2,
\mathbf{L}_{-F}\phi_3
\right)
\widetilde{\rho}_{\tilde{a}}\!\left(
\widetilde{\mathbf{L}}_{-\widetilde B}\widetilde{\phi}_1,
\widetilde{\mathbf{L}}_{-\widetilde D}\widetilde{\phi}_2,
\widetilde{\mathbf{L}}_{-\widetilde F}\widetilde{\phi}_3
\right).
\fe
Here $a= A+C+E~\pmod 2, \tilde{a}=\tilde{A} +\tilde{C}+\tilde{E} \pmod 2$. It satisfy $a+p_1+p_2+p_3 = \tilde{a}+\tilde{p_1}+\tilde{p_2}+\tilde{p_3} \pmod 2$.
In this step, the sign comming from the reordering of holomorphic primaries and anti-holomorphic primaries will cancel out because the same factor appears in both matrix element.

Using the factorialization properies, the genus 2 partition function simplifies into
\ie
Z_{g=2}
={}&
\sum_{\mathcal{O}_1,\mathcal{O}_2,\mathcal{O}_3}
\sum_{\substack{
|A|=|B|,\ |C|=|D|,\ |E|=|F|\\
|\widetilde A|=|\widetilde B|,\ 
|\widetilde C|=|\widetilde D|,\ 
|\widetilde E|=|\widetilde F|
}}
\prod_{i=1}^{3}
q_i^{h_i}\widetilde q_i^{\widetilde h_i} \eta_i^{p_i} \widetilde \eta_i^{\tilde{p}_i}
\times
q_1^{|A|}\eta_1^A
q_2^{|C|}\eta_2^C
q_3^{|E|}\eta_3^E
\widetilde q_1^{|\widetilde A|}
\widetilde\eta_1^{\widetilde A}
\widetilde q_2^{|\widetilde C|}
\widetilde\eta_2^{\widetilde C}
\widetilde q_3^{|\widetilde E|}
\widetilde\eta_3^{\widetilde E}
\\
&\times
(-1)^{a+p_1+p_2+p_3+K +\bar{K}}
\left(C_{123}^{(a)}\right)^2 \mathbf{G}_{h_1}^{AB}
\mathbf{G}_{h_2}^{CD}
\mathbf{G}_{h_3}^{EF}
\widetilde{\mathbf{G}}_{\widetilde h_1}^{\widetilde A\widetilde B}
\widetilde{\mathbf{G}}_{\widetilde h_2}^{\widetilde C\widetilde D}
\widetilde{\mathbf{G}}_{\widetilde h_3}^{\widetilde E\widetilde F}
\\
&\quad\times
\rho_a\!\left(
\mathbf{L}_{-A}\phi_1,
\mathbf{L}_{-C}\phi_2,
\mathbf{L}_{-E}\phi_3
\right)
\widetilde{\rho}_{\tilde{a}}\!\left(
\widetilde{\mathbf{L}}_{-\widetilde A}\widetilde{\phi}_1,
\widetilde{\mathbf{L}}_{-\widetilde C}\widetilde{\phi}_2,
\widetilde{\mathbf{L}}_{-\widetilde E}\widetilde{\phi}_3
\right)
\\
&\quad\times
\rho_a\!\left(
\mathbf{L}_{-B}\phi_1,
\mathbf{L}_{-D}\phi_2,
\mathbf{L}_{-F}\phi_3
\right)
\widetilde{\rho}_{\tilde{a}}\!\left(
\widetilde{\mathbf{L}}_{-\widetilde B}\widetilde{\phi}_1,
\widetilde{\mathbf{L}}_{-\widetilde D}\widetilde{\phi}_2,
\widetilde{\mathbf{L}}_{-\widetilde F}\widetilde{\phi}_3
\right).
\fe
Here
\ie
K
\equiv{}&
(A+p_1)(C+p_2)
+(A+p_1)(E+p_3)
+(C+p_2)(E+p_3)
\pmod{2},
\\
\bar{K}
\equiv{}&
(\widetilde A+\widetilde p_1)
(\widetilde C+\widetilde p_2)
+
(\widetilde A+\widetilde p_1)
(\widetilde E+\widetilde p_3)
+
(\widetilde C+\widetilde p_2)
(\widetilde E+\widetilde p_3)
\pmod{2}.
\fe
Define the NS superconformal block by
\ie \label{NSblockThetaDefinition}
\mathbf{F}^{\Theta}_a(h_1, h_2, h_3, c)= & \sum_{\substack{|A|=|B|,|C|=|D|,|E|=|F|,\\ (-1)^{A+C+E}=(-1)^a}} q_1^{|A|}\eta_1^{A+p_1} q_2^{|C|} \eta_2^{C+p_2} q_3^{|E|} \eta_3^{E+p_3} \mathbf{G}_{h_1}^{A B} \mathbf{G}_{h_2}^{C D} \mathbf{G}_{h_3}^{E F} \\
& \times (-1)^{K}\rho_a(\mathbf{L}_{-A} \phi_1, \mathbf{L}_{-C} \phi_2, \mathbf{L}_{-E} \phi_3) \rho_a(\mathbf{L}_{-B} \phi_1, \mathbf{L}_{-D} \phi_2, \mathbf{L}_{-F} \phi_3) ,
\fe
and analogously the anti-holomorphic block $\widetilde{F}_{\tilde{a}}$. The genus 2 partition function then reduces to
\ie
Z_{g=2} = \sum_{\mathcal{O}_1,\mathcal{O}_2,\mathcal{O}_3} \sum_{a=0,1}  \left(C^{(a)}_{123} \right)^2(-1)^{a+p_1+p_2+p_3} \prod_{i=1}^{3}
q_i^{h_i}\widetilde q_i^{\widetilde h_i} F_a(h_1, h_2, h_3, c) \widetilde{F}_{\tilde{a}} (\tilde{h}_1,\tilde{h}_2, \tilde{h}_3, \tilde{c}).
\fe
where $a+p_1+p_2+p_3 = \tilde{a}+\tilde{p_1}+\tilde{p_2}+\tilde{p_3} \pmod 2$.

\subsubsection*{All-NS Block Decomposition in Glass Channel}
We can also compute the genus 2 partition function in the glass channel. It is given by the following prescription
\ie
Z^{Gl}_{g=2}
={}&
\sum_{\mathcal{O}_1,\mathcal{O}_2,\mathcal{O}_3}
\sum_{\substack{
|A|=|B|,\ |C|=|D|,\ |E|=|F|\\
|\widetilde A|=|\widetilde B|,\ 
|\widetilde C|=|\widetilde D|,\ 
|\widetilde E|=|\widetilde F|
}}
\prod_{i=1}^{3}
q_i^{h_i}\widetilde q_i^{\widetilde h_i} \eta_i^{p_i} \widetilde \eta_i^{\tilde{p}_i}
\times
q_1^{|A|}\eta_1^A
q_2^{|C|}\eta_2^C
q_3^{|E|}\eta_3^E
\widetilde q_1^{|\widetilde A|}
\widetilde\eta_1^{\widetilde A}
\widetilde q_2^{|\widetilde C|}
\widetilde\eta_2^{\widetilde C}
\widetilde q_3^{|\widetilde E|}
\widetilde\eta_3^{\widetilde E}
\\
&\times
(-1)^{
(E+\widetilde E+p_3+\widetilde p_3)
\left[
(A+\widetilde A+p_1+\widetilde p_1)
+
(C+\widetilde C+p_2+\widetilde p_2)
\right]
}
\\
&\times
(\mathbb{G}_1)^{A\widetilde A;B\widetilde B}
(\mathbb{G}_2)^{C\widetilde C;D\widetilde D}
(\mathbb{G}_3)^{E\widetilde E;F\widetilde F}
\\
&\times
\left\langle\!\left\langle
\mathbf{L}_{-A} \widetilde{\mathbf{L}}_{-\widetilde A}\mathcal{O}_1
\middle|
\mathbf{L}_{-E} \widetilde{\mathbf{L}}_{-\widetilde E}\mathcal{O}_3(1,1)
\middle|
\mathbf{L}_{-B} \widetilde{\mathbf{L}}_{-\widetilde B}\mathcal{O}_1
\right\rangle\right.
\times
\left\langle\!\left\langle
\mathbf{L}_{-C} \widetilde{\mathbf{L}}_{-\widetilde C}\mathcal{O}_2
\middle|
\mathbf{L}_{-F} \widetilde{\mathbf{L}}_{-\widetilde F}\mathcal{O}_3(1,1)
\middle|
\mathbf{L}_{-D} \widetilde{\mathbf{L}}_{-\widetilde D}\mathcal{O}_2
\right\rangle\right.
\fe
The gram matrix factorize in the same way as \eqref{eq: Gram_Factorialization}. The three point function would factorize as
\ie
&\left\langle\!\left\langle
\mathbf{L}_{-A}
\widetilde{\mathbf{L}}_{-\widetilde A}\mathcal{O}_1
\middle|
\mathbf{L}_{-E}
\widetilde{\mathbf{L}}_{-\widetilde E}\mathcal{O}_3(1,1)
\middle|
\mathbf{L}_{-B}
\widetilde{\mathbf{L}}_{-\widetilde B}\mathcal{O}_1
\right\rangle\right.
\\
&\quad\times
\left\langle\!\left\langle
\mathbf{L}_{-C}
\widetilde{\mathbf{L}}_{-\widetilde C}\mathcal{O}_2
\middle|
\mathbf{L}_{-F}
\widetilde{\mathbf{L}}_{-\widetilde F}\mathcal{O}_3(1,1)
\middle|
\mathbf{L}_{-D}
\widetilde{\mathbf{L}}_{-\widetilde D}\mathcal{O}_2
\right\rangle\right.
\\
={}&
(-1)^{
p_1\widetilde p_1+p_2\widetilde p_2+p_3\widetilde p_3
+(A+p_1+E+p_3)(\widetilde A+\widetilde p_1+\widetilde E+\widetilde p_3)
+(C+p_2+E+p_3)(\widetilde C+\widetilde p_2+\widetilde E+\widetilde p_3)
}
D_1D_2D_3
\left(C_{131}^{(a)}C_{232}^{(a)}\right)
\\
&\quad\times
\rho_a\!\left(
\mathbf{L}_{-A}\phi_1,
\mathbf{L}_{-E}\phi_3,
\mathbf{L}_{-B}\phi_1
\right)
\widetilde{\rho}_{\tilde{a}}\!\left(
\widetilde{\mathbf{L}}_{-\widetilde A}\widetilde{\phi}_1,
\widetilde{\mathbf{L}}_{-\widetilde E}\widetilde{\phi}_3,
\widetilde{\mathbf{L}}_{-\widetilde B}\widetilde{\phi}_1
\right)
\\
&\quad\times
\rho_a\!\left(
\mathbf{L}_{-C}\phi_2,
\mathbf{L}_{-F}\phi_3,
\mathbf{L}_{-D}\phi_2
\right)
\widetilde{\rho}_{\tilde{a}}\!\left(
\widetilde{\mathbf{L}}_{-\widetilde C}\widetilde{\phi}_2,
\widetilde{\mathbf{L}}_{-\widetilde F}\widetilde{\phi}_3,
\widetilde{\mathbf{L}}_{-\widetilde D}\widetilde{\phi}_2
\right).
\fe
where $a\equiv E\pmod 2$ and  $\tilde{a}= \widetilde{E} \pmod 2$.

Using the factorialization properties, the glass channel genus 2 partition function can be expressed into
\[
\begin{aligned}
Z_{g=2}^{\mathrm{Gl}}
={}&
\sum_{\mathcal O_1,\mathcal O_2,\mathcal O_3}
\sum_{\substack{
|A|=|B|,\ |C|=|D|,\ |E|=|F|\\
|\widetilde A|=|\widetilde B|,\ 
|\widetilde C|=|\widetilde D|,\ 
|\widetilde E|=|\widetilde F|
}}
\prod_{i=1}^{3}
q_i^{h_i}\widetilde q_i^{\widetilde h_i} \eta_i^{p_i} \widetilde{\eta}_i^{\tilde{p}_i}
q_1^{|A|}\eta_1^A
q_2^{|C|}\eta_2^C
q_3^{|E|}\eta_3^E
\widetilde q_1^{|\widetilde A|}\widetilde\eta_1^{\widetilde A}
\widetilde q_2^{|\widetilde C|}\widetilde\eta_2^{\widetilde C}
\widetilde q_3^{|\widetilde E|}\widetilde\eta_3^{\widetilde E}
\\
&\times
C_{131}^{(a)}C_{232}^{(a)} ~(-1)^{a+p_3+K_{\mathrm{Gl}}+\widetilde K_{\mathrm{Gl}}}
~\mathbf G_{h_1}^{AB}
\mathbf G_{h_2}^{CD}
\mathbf G_{h_3}^{EF}
\widetilde{\mathbf G}_{\widetilde h_1}^{\widetilde A\widetilde B}
\widetilde{\mathbf G}_{\widetilde h_2}^{\widetilde C\widetilde D}
\widetilde{\mathbf G}_{\widetilde h_3}^{\widetilde E\widetilde F}
\\
&\times
\rho_a\!\left(
\mathbf L_{-A}\phi_1,
\mathbf L_{-E}\phi_3,
\mathbf L_{-B}\phi_1
\right)
\rho_a\!\left(
\mathbf L_{-C}\phi_2,
\mathbf L_{-F}\phi_3,
\mathbf L_{-D}\phi_2
\right)
\\
&\times
\widetilde\rho_a\!\left(
\widetilde{\mathbf L}_{-\widetilde A}\widetilde\phi_1,
\widetilde{\mathbf L}_{-\widetilde E}\widetilde\phi_3,
\widetilde{\mathbf L}_{-\widetilde B}\widetilde\phi_1
\right)
\widetilde\rho_a\!\left(
\widetilde{\mathbf L}_{-\widetilde C}\widetilde\phi_2,
\widetilde{\mathbf L}_{-\widetilde F}\widetilde\phi_3,
\widetilde{\mathbf L}_{-\widetilde D}\widetilde\phi_2
\right).
\end{aligned}
\]
where 
\ie
K_{\mathrm{Gl}}
=a\bigl(A+p_1+C+p_2\bigr),\quad
\widetilde K_{\mathrm{Gl}} =a\bigl(\widetilde A+\widetilde p_1
+\widetilde C+\widetilde p_2\bigr).
\fe
Define the NS superconformal block in glass channel as
\ie
 \label{NSblockGlassesDefinition}
\mathbf{F}^{\mathrm{Gl}}_a(h_1, h_2, h_3, c)= & \sum_{|A|=|B|,|C|=|D|,|E|=|F|, (-1)^{E}=(-1)^a} q_1^{|A|}\eta_1^{A+p_1} q_2^{|C|} \eta_2^{C+p_2} q_3^{|E|} \eta_3^{E+p_3} \mathbf{G}_{h_1}^{A B} \mathbf{G}_{h_2}^{C D} \mathbf{G}_{h_3}^{E F} \\
& \times (-1)^{K_{\mathrm{Gl}}}\rho_a(\mathbf{L}_{-A} \phi_1, \mathbf{L}_{-E} \phi_3, \mathbf{L}_{-B} \phi_1) \rho_a(\mathbf{L}_{-C} \phi_2, \mathbf{L}_{-F} \phi_3, \mathbf{L}_{-D} \phi_2).
\fe
and analogously the anti-holomorphic block $\widetilde{F}^{\mathrm{Gl}}_a$. The genus 2 partition function then reduces to
\ie
Z_{g=2} = \sum_{\mathcal{O}_1,\mathcal{O}_2,\mathcal{O}_3} \sum_{a=0,1}  \left(C^{(a)}_{131} C^{(a)}_{232}  \right) (-1)^{a+p_1} \prod_{i=1}^{3}
q_i^{h_i}\widetilde q_i^{\widetilde h_i} F^{\mathrm{Gl}}_a(h_1, h_2, h_3, c) \widetilde{F}^{\mathrm{Gl}}_{\tilde{a}} (\tilde{h}_1,\tilde{h}_2, \tilde{h}_3, \tilde{c}).
\fe
where $a+p_1= \tilde{a}+\tilde{p}_1 \pmod 2$.


\section{NS superconformal blocks}


\subsection{c-Recursion for NS Block}

We will keep using the theta channel as an example, with the plumbing data denoted by $(q_i, \eta_i)$. The all-NS conformal block is given by
\ie \label{NSblockRecursionDefinition}
\mathbf{F}_a(h_1, h_2, h_3, c)= &\sum_{\substack{
|A|=|B|,\ |C|=|D|,\ |E|=|F|\\
(-1)^{A+C+E}=(-1)^a
}}q_1^{|A|}\eta_1^{A+p_1} q_2^{|C|} \eta_2^{C+p_2} q_3^{|E|} \eta_3^{E+p_3} \mathbf{G}_{h_1}^{A B} \mathbf{G}_{h_2}^{C D} \mathbf{G}_{h_3}^{E F} \\
& \times (-1)^{(A+p_1)(C+p_2)+(A+p_1)(E+p_3)+(C+p_2)(E+p_3)} \\
&\rho_a(\mathbf{L}_{-A} \phi_1, \mathbf{L}_{-C} \phi_2, \mathbf{L}_{-E} \phi_3) \rho_a(\mathbf{L}_{-B} \phi_1, \mathbf{L}_{-D} \phi_2, \mathbf{L}_{-F} \phi_3) ,
\fe
where $\mathbf{G}$ denotes the inverse of the super-Gram matrix $\langle h|\mathbf{L}_{-A}^\dagger\mathbf{L}_{-B}|h\rangle$. $a$ takes value from $0, 1$. The sign $(-1)^{(A+p_1)(C+p_2)+(A+p_1)(E+p_3)+(C+p_2)(E+p_3)}$ comes from moving the operator orders. The condition $(-1)^{A+C+E}=(-1)^a$ ensures that only total-integer level or total-half-integer level intermediate states enters the 3-point function.

The NS supermultiplet admits a level-$rs/2$ null state when
\ie
h_{rs}(c) = \frac{Q^2-(rb+sb^{-1})^2}{8}, \quad Q=b+b^{-1}\quad (r\geq 2,s\geq 1, \ r+s\in 2\mathbb{Z}),
\fe
where $b$ is defined by $c= \frac{3}{2}+3(b+b^{-1})^2$. Conversely, for fixed $h$, the degeneration happens when the central charge is
\ie
c_{r,s}(h) = \frac{15}{2}+3(x_{rs}+x_{rs}^{-1}),\quad b_{rs} = \sqrt{x_{rs}},\quad x_{rs} = \frac{r s-1+4 h+\sqrt{(r-s)^2+8(r s-1) h+16 h^2}}{1-r^2}.
\fe
At $c=c_{r,s}$, the Gram matrix becomes degenerate, which indicate a pole in the conformal block. For generic $h_i$, the pole is a simple pole. It suggests the following structure for the all-NS conformal block
\ie
\mathbf{F}_a(h_1, h_2, h_3, c)= \mathbf{F}^{\text{global}}_a(h_1,h_2,h_3,c)+\sum_{i=1}^3 \sum_{r\geq 2,s\geq 1} \frac{\mathbf{W}^{(i)}_{a,rs}(h_1,h_2,h_3)}{c-c_{r,s}(h_i)}
\fe
where $\mathbf{F}^{\text{global}}_a(h_1,h_2,h_3,c)$ is regular on the entire $c$ plane. Follow from the representation analysis as for the Virasoro block, we can determine that $\mathbf{F}^{\text{global}}_a(h_1,h_2,h_3,c)$ is also regular at $c\rightarrow \infty$. It determines that $ \mathbf{F}^{\text{global}}_a(h_1,h_2,h_3,c)$ must be independent of $c$. Therefore, we can access the $\mathbf{F}^{\text{global}}_a$ by taking the limit
\ie
\mathbf{F}^{\text{global}}_a(h_1,h_2,h_3) = \lim_{c\rightarrow \infty} \mathbf{F}_a(h_1, h_2, h_3, c).
\fe
It would be referred to as the large-$c$ seed term. 

The subalgebra of $\mathcal{N}=1$ SCA that's independent of $c$ is given by the $\mathfrak{osp}(1|2)$ algebra generated by $\{L_{0,\pm 1}, G_{\pm 1/2}\}$. We thus define $\hat{A}$ as the number of non-$L_{-1}, G_{-1/2}$ elements in $\mathbf{L}_{-A}$. The same $c$-power counting argument works also in the SUSY case: If we construct a orthonormal basis $\mathcal{L}_{-A}$ by applying Gram-Schmidt on $\mathbf{L}_{-A}$, then in the large $c$ limit the 3-point function $\rho(\mathcal{L}_{-A}\phi_1,\mathcal{L}_{-B}\phi_2, \mathcal{L}_{-C}\phi_3)$ factorizes into the large-$c$ vacuum block and the $\mathfrak{osp}(1|2)$ block,
\ie
\lim_{c\rightarrow \infty}\mathbf{F}_a(h_1, h_2, h_3,c) &= \sum_{\substack{a_1,a_2,a_3\in\{0,1\}\\a_1+a_2+a_3\equiv a\ (\mathrm{mod}\ 2)}} \mathbf{F}_0(0,0,0,\infty;\eta_i\rightarrow (-1)^{a_i+p_i}\eta_i) \times \mathbf{F}_{a_1a_2a_3}^{\mathfrak{osp}(1|2)}(h_1,h_2,h_3),
\fe
where the $\mathfrak{osp}(1|2)$ block is defined for $a_i\in\{0,1\}$ and $a_1+a_2+a_3\equiv a\ (\mathrm{mod}\ 2)$ by
\ie
\mathbf{F}^{\mathfrak{osp}(1|2)}_{a_1a_2a_3}(h_1, h_2, h_3) &= \sum_{k_1, k_2, k_3} q_1^{k_1+\frac{1}{2}a_1}\eta_1^{p_1+a_1}q_2^{k_2+\frac{1}{2}a_2}\eta_2^{p_2+a_2}q_3^{k_3+\frac{1}{2}a_3}\eta_3^{p_3+a_3} \\&\quad \times (-1)^{(a_1+p_1)(a_2+p_2)+(a_1+p_1)(a_3+p_3)+(a_2+p_2)(a_3+p_3)}\\
&\quad \times \left[\prod_{i=1}^3 \frac{1}{\|L_{-1}^{k_i}G_{-\frac{1}{2}}^{a_i}\phi_i\|^2}\right] \rho_a(L_{-1}^{k_1}G_{-\frac{1}{2}}^{a_1}\phi_1,L_{-1}^{k_2}G_{-\frac{1}{2}}^{a_2}\phi_2,L_{-1}^{k_3}G_{-\frac{1}{2}}^{a_3}\phi_3)^2
\fe
Importantly, the large-$c$ limit of the odd vacuum block is 0, due to the fact that one can only contract $G_r,G_{-r}$ pairwise. Due to the Grassmann parity of $G_{-1/2}^{a}\phi$, for each fixed $a$ one must sum over the corresponding 4 configurations of global $\mathfrak{sl}(2)$ primaries independently; the two values of $a$ give 8 configurations in total. The sign $\eta_i$ in the corresponding global block must be shifted. Similar to the Virasoro case, the vacuum block can also be directly computed in Schottky parameters. We will focus on the global $\mathfrak{osp}(1|2)$ blocks in the next section.



\subsection{Global $\mathfrak{osp}(1|2)$ blocks}

In this section, we develop the method of computing arbitrary Global $\mathfrak{osp}(1|2)$ blocks. We take the following basis for states
\ie
L_{-1}^k G_{-\frac{1}{2}}^a \phi,\quad k\in \mathbb{N},\ a\in \{0,1\}.
\fe
with $\phi$ normalized such that $\|\phi\|^2=1$. The nontrivial inner products are
\ie
\|L_{-1}^k\phi\|^2 = k!(2h)_k,\quad \|L_{-1}^k G_{-\frac{1}{2}}\phi\|^2 = k! (2h)_{k+1},
\fe
where $(x)_n$ is the rising Pochhammer symbol.

For the genus-2 block in the theta channel, fix $a_i\in\{0,1\}$ and set $a\equiv a_1+a_2+a_3\ (\mathrm{mod}\ 2)$. We then have
\ie
\mathbf{F}^{\mathfrak{osp}(1|2)}_{a_1a_2a_3}(h_1, h_2, h_3) &= \sum_{k_1, k_2, k_3} q_1^{k_1+\frac{1}{2}a_1}\eta_1^{a_1+p_1}q_2^{k_2+\frac{1}{2}a_2}\eta_2^{a_2+p_2}q_3^{k_3+\frac{1}{2}a_3}\eta_3^{a_3+p_3} \\&\quad \times (-1)^{(a_1+p_1)(a_2+p_2)+(a_1+p_1)(a_3+p_3)+(a_2+p_2)(a_3+p_3)}\\
&\quad \times \left[\prod_{i=1}^3 \frac{1}{k_i!(2h_i)_{k_i+a_i}}\right] \rho_a(L_{-1}^{k_1}G_{-\frac{1}{2}}^{a_1}\phi_1,L_{-1}^{k_2}G_{-\frac{1}{2}}^{a_2}\phi_2,L_{-1}^{k_3}G_{-\frac{1}{2}}^{a_3}\phi_3)^2
\fe
The remaining object to compute is simply $\rho_a(L_{-1}^{k_1}G_{-\frac{1}{2}}^{a_1}\phi_1,L_{-1}^{k_2}G_{-\frac{1}{2}}^{a_2}\phi_2,L_{-1}^{k_3}G_{-\frac{1}{2}}^{a_3}\phi_3)$, which can be computed directly since $G^a_{-\frac{1}{2}}\phi$ is a $\mathfrak{sl}(2)$ primary and one can use the result for $\mathfrak{sl}(2)$ 3-point functions.  We keep the normalization 
\ie
\rho_0(\phi_1,\phi_2,\phi_3)=1,\quad \rho_1(\phi_1,G_{-1/2}\phi_2,\phi_3)=1.
\fe
Superconformal Wald identity fix the rest to be 
\[
\begin{array}{c|c|c}
(a_1a_2a_3)
&a&
\rho_a\left(
G_{-\frac12}^{a_1}\phi_1,
G_{-\frac12}^{a_2}\phi_2,
G_{-\frac12}^{a_3}\phi_3
\right)
\\ \hline
000&0&1\\
110&0&(-1)^{p_1}(h_1+h_2-h_3)\\
101&0&(-1)^{p_1+p_2}(h_1-h_2+h_3)\\
011&0&(-1)^{p_2}(h_1-h_2-h_3)\\ \hline
100&1&(-1)^{p_1}\\
010&1&1\\
001&1&(-1)^{p_2+1}\\
111&1&(-1)^{p_1+p_2+1}
\left(h_1+h_2+h_3-\frac12\right)
\end{array}
\]
Defining 
\ie
s_{km}(h_1,h_2,h_3)
 &=
 \sum_{p=0}^{\min(k,m)}
 \frac{k!}{p!(k-p)!}
 (2h_3+m-1)^{(p)}m^{(p)}
 \\
 &\qquad\quad\times
 (h_3+h_2-h_1)_{m-p}
 (h_1+h_2-h_3+p-m)_{k-p},
 \label{eq:ccy-s-formula}
\fe
where $(x)^{(n)}$ is the falling Pochhammer symbol. Using the formula for $\mathfrak{sl}(2)$ 3-point functions in \cite{Cho-Collier-Yin}, one then arrive at
\ie
 &\rho_a\!\left(
 L_{-1}^{k_1}G_{-1/2}^{a_1}\phi_1,
 L_{-1}^{k_2}G_{-1/2}^{a_2}\phi_2,
 L_{-1}^{k_3}G_{-1/2}^{a_3}\phi_3
 \right)
 \\
 &\quad =
 \rho_a\!\left(
 G_{-1/2}^{a_1}\phi_1,
 G_{-1/2}^{a_2}\phi_2,
 G_{-1/2}^{a_3}\phi_3
 \right) s_{k_1k_3}\!\left(
 h_1+\frac{a_1}{2},
 h_2+\frac{a_2}{2},
 h_3+\frac{a_3}{2}
 \right)
 \\
 &\qquad\quad\times
 \left(
 h_1+\frac{a_1}{2}+k_1
 -h_2-\frac{a_2}{2}-k_2+1
 -h_3-\frac{a_3}{2}-k_3
 \right)_{k_2},
\fe
and hence
\ie
\mathbf{F}^{\mathfrak{osp}(1|2)}_{a_1a_2a_3}(h_1, h_2, h_3) &= \sum_{k_1, k_2, k_3} q_1^{k_1+\frac{1}{2}a_1}\eta_1^{a_1+p_1}q_2^{k_2+\frac{1}{2}a_2}\eta_2^{a_2+p_2}q_3^{k_3+\frac{1}{2}a_3}\eta_3^{a_3+p_3} \\&\quad \times (-1)^{(a_1+p_1)(a_2+p_2)+(a_1+p_1)(a_3+p_3)+(a_2+p_2)(a_3+p_3)}
\left[\prod_{i=1}^3 \frac{1}{k_i!(2h_i)_{k_i+a_i}}\right]\\
&\quad \times \Bigg[\rho_a\!\left(
 G_{-1/2}^{a_1}\phi_1,
 G_{-1/2}^{a_2}\phi_2,
 G_{-1/2}^{a_3}\phi_3
 \right) s_{k_1k_3}\!\left(
 h_1+\frac{a_1}{2},
 h_2+\frac{a_2}{2},
 h_3+\frac{a_3}{2}
 \right)
 \\
 &\quad\qquad \times
 \left(
 h_1+\frac{a_1}{2}+k_1
 -h_2-\frac{a_2}{2}-k_2+1
 -h_3-\frac{a_3}{2}-k_3
 \right)_{k_2}\Bigg]^2.
\fe
Using this formula, one can compute the $\mathfrak{osp}(1|2)$ block directly as a $q$-expansion. This concludes our discussion for the all-NS $c$-recursion relation.




\subsection{Residue in the NS Block c-Recursion}

At a pole $c_{rs}$ on the complex $c$-plane, the residues are contributed by the descendents of this null state, which we will call $\chi_{rs}=\sum_{M} \chi_{rs}^M \mathbf{L}_{-M}\phi$. 
The pole of $\mathbf{F}_a$ at $c=c_{rs}(h_1)$ is given by
\ie
\mathbf{W}^{(1)}_{a, rs} &= -\frac{\partial c_{rs}(h_1)}{\partial h_1} \lim_{h_1\rightarrow h_{r,s}} \frac{h_1-h_{rs}}{\|\sum_{M} \chi_{rs}^M \mathbf{L}_{-M}\phi_1 \|^2} \\
&\times  \sum_{\substack{
|A|=|B|,\ |C|=|D|,\ |E|=|F|\\
(-1)^{A+C+E}=(-1)^{a+rs}
}} q_1^{|A|+\frac{rs}{2}}\eta_1^{A+p_1+rs} q_2^{|C|} \eta_2^{C+p_2} q_3^{|E|} \eta_3^{E+p_3} \mathbf{G}_{h_1+\frac{rs}{2}}^{A B} \mathbf{G}_{h_2}^{C D} \mathbf{G}_{h_3}^{E F} \\
&\quad \quad \quad (-1)^{(A+p_1+rs)(C+p_2)+(A+p_1+rs)(E+p_3)+(C+p_2)(E+p_3)}\\
& \quad \quad \quad\rho_a(\mathbf{L}_{-A} \chi_{rs}, \mathbf{L}_{-C} \phi_2, \mathbf{L}_{-E} \phi_3) \rho_a(\mathbf{L}_{-B} \chi_{rs}, \mathbf{L}_{-D} \phi_2, \mathbf{L}_{-F} \phi_3).
\fe
Normalizing $\chi_{rs}$ by taking the coefficient of $G_{-1/2}^{rs}\phi$ to be 1, the inverse residue of the null norm is given by
\ie
\mathbf{A}_{r s}^c(h)=\lim_{h_1\rightarrow h_{r,s}} \frac{h_1-h_{rs}}{\|\sum_{M} \chi_{rs}^M \mathbf{L}_{-M}\phi_1 \|^2}=\frac{1}{2} \prod_{\substack{p=1-r, \ldots, r \\ q=1-s, \ldots, s \\ p+q \in 2\mathbb{Z} \\(p, q) \neq(0,0),(r, s)}}\left[\frac{p b_{r s}+q b_{r s}^{-1}}{\sqrt{2}}\right]^{-1}
\fe
The 3-point functions can be factorized as
\ie
\rho_a(\mathbf{L}_{-A}\chi_{rs} , \mathbf{L}_{-B}\phi_2,  \mathbf{L}_{-C}\phi_3) = \begin{cases} \rho_{a}(\mathbf{L}_{-A}\phi_{h_{rs}+\frac{rs}{2}} , \mathbf{L}_{-B}\phi_2,  \mathbf{L}_{-C}\phi_3) \mathbf{P}_{rs}^a(h_3, h_2),\quad rs\in 2\mathbb{Z}, \\ 
(-1)^{p_1+A}\rho_{1-a}(\mathbf{L}_{-A}\phi_{h_{rs}+\frac{rs}{2}} , \mathbf{L}_{-B}\phi_2,  \mathbf{L}_{-C}\phi_3) \mathbf{P}_{rs}^{a}(h_3, h_2),\quad rs\in 2\mathbb{Z}+1,
\end{cases}
\fe
where the two fusion polynomials are given by
\ie
\mathbf{P}_{rs}^0(h_i, h_j;c)&=\prod_{\substack{p=1-r, 1-r+2, \ldots, r-1 \\ q=1-s, 1-s+2, \ldots, s-1 \\ p+q-(r+s) \in 4 \mathbb{Z}+2}} \frac{\lambda_i-\lambda_j+p b+q b^{-1}}{2 \sqrt{2}} \frac{\lambda_i+\lambda_j+p b+q b^{-1}}{2 \sqrt{2}},\\
\mathbf{P}_{rs}^1(h_i, h_j;c) &= \prod_{\substack{p=1-r, 1-r+2, \ldots, r-1 \\ q=1-s, 1-s+2, \ldots, s-1 \\ p+q-(r+s) \in 4 \mathbb{Z}}} \frac{\lambda_i-\lambda_j+p b+q b^{-1}}{2 \sqrt{2}} \frac{\lambda_i+\lambda_j+p b+q b^{-1}}{2 \sqrt{2}}.
\fe
Here $\lambda_i$ is defined by
\ie
h_i=\frac{(b+b^{-1})^2-\lambda_i^2}{8}.
\fe
Plugging these results in the pole $\mathbf{W}^{(1)}_{a,rs}$, one then have
\ie
\mathbf{W}^{(1)}_{a,rs} &= -\frac{\partial c_{rs}(h_1)}{\partial h_1} \mathbf{A}_{rs}^c(h) q_1^{rs/2}\eta_1^{rs} [\mathbf{P}_{rs}^a(h_3, h_2)]^2\\
&\quad \times \begin{cases}
\mathbf{F}_{a}(h_1 \rightarrow h_1+\tfrac{rs}{2}, c\rightarrow c_{rs}),\quad rs\in 2\mathbb{Z}, & rs\in 2\mathbb{Z},\\
\mathbf{F}_{1-a}(h_1 \rightarrow h_1+\tfrac{rs}{2}, c\rightarrow c_{rs},\eta_{2,3}\rightarrow -\eta_{2,3}),& rs\in 2\mathbb{Z}+1,
\end{cases}
\fe
where the additional shift $\eta_{2,3}\rightarrow - \eta_{2,3}$ comes from the additional cross sign factor $(-1)^{rs(C+p_2+E+p_3)}$.

Similar logic can be applied to compute the residue at $c=c_{r,s}(h_2)$ and $c=c_{r,s}(h_3)$. For that computation, one need to be careful about the sign in the factorialization of three point function. For a null state in the second slot, one have
\ie
&\rho_a\!\left(
\mathbf L_{-A}\phi_1,
\mathbf L_{-C}\chi_{rs},
\mathbf L_{-E}\phi_3
\right)
\\
&\quad=
\begin{cases}
\mathbf P_{rs}^{a}(h_3,h_1;c_{rs})\rho_a\!\left(
\mathbf L_{-A}\phi_1,
\mathbf L_{-C}\phi_{h_{rs}+\frac{rs}{2}},
\mathbf L_{-E}\phi_3
\right),
& rs\in2\mathbb Z,
\\[0.4em]
\mathbf P_{rs}^{a}(h_3,h_1;c_{rs})\rho_{1-a}\!\left(
\mathbf L_{-A}\phi_1,
\mathbf L_{-C}\phi_{h_{rs}+\frac{rs}{2}},
\mathbf L_{-E}\phi_3
\right),
& rs\in2\mathbb Z+1.
\end{cases}
\fe 
For a null state in the third slot,
\ie
&\rho_a\!\left(
\mathbf L_{-A}\phi_1,
\mathbf L_{-C}\phi_2,
\mathbf L_{-E}\chi_{rs}
\right)
\\
&\quad=
\begin{cases}
\mathbf P_{rs}^{a}(h_1,h_2;c_{rs})\rho_a\!\left(
\mathbf L_{-A}\phi_1,
\mathbf L_{-C}\phi_2,
\mathbf L_{-E}\phi_{h_{rs}+\frac{rs}{2}}
\right),
& rs\in2\mathbb Z,
\\[0.4em]
(-1)^{1+p_2}
\mathbf P_{rs}^{a}(h_1,h_2;c_{rs})\rho_{1-a}\!\left(
\mathbf L_{-A}\phi_1,
\mathbf L_{-C}\phi_2,
\mathbf L_{-E}\phi_{h_{rs}+\frac{rs}{2}}
\right),
& rs\in2\mathbb Z+1.
\end{cases}
\fe
The additional signs will cancel out in the theta channel, giving

\ie
\mathbf{W}^{(2)}_{a,rs} &= -\frac{\partial c_{rs}(h_2)}{\partial h_2} \mathbf{A}_{rs}^c(h) q_2^{rs/2}\eta_2^{rs} [\mathbf{P}_{rs}^a(h_3, h_1)]^2\\
&\quad \times \begin{cases}
\mathbf{F}_{a}(h_2 \rightarrow h_2+\tfrac{rs}{2}, c\rightarrow c_{rs}),\quad rs\in 2\mathbb{Z}, & rs\in 2\mathbb{Z},\\
\mathbf{F}_{1-a}(h_2 \rightarrow h_2+\tfrac{rs}{2}, c\rightarrow c_{rs},\eta_{3,1}\rightarrow -\eta_{3,1}),& rs\in 2\mathbb{Z}+1,
\end{cases}
\fe

\ie
\mathbf{W}^{(3)}_{a,rs} &= -\frac{\partial c_{rs}(h_3)}{\partial h_3} \mathbf{A}_{rs}^c(h) q_3^{rs/2}\eta_3^{rs} [\mathbf{P}_{rs}^a(h_1, h_2)]^2\\
&\quad \times \begin{cases}
\mathbf{F}_{a}(h_3 \rightarrow h_3+\tfrac{rs}{2}, c\rightarrow c_{rs}),\quad rs\in 2\mathbb{Z}, & rs\in 2\mathbb{Z},\\
\mathbf{F}_{1-a}(h_3 \rightarrow h_3+\tfrac{rs}{2}, c\rightarrow c_{rs},\eta_{1,2}\rightarrow -\eta_{1,2}),& rs\in 2\mathbb{Z}+1.
\end{cases}
\fe




\section{The Ramond blocks}

Let us now consider the case with intermediate Ramond supermultiplets. The $h$-recursion relation of sphere 4-point Ramond block is derived in \cite{0810.1203, 1012.2974}, and that of torus 1-point Ramond block is derived in \cite{1207.5740}. In this section, we aim for $c$-recursion of arbitrary Ramond blocks.

A Ramond multiplet has 2 ground states $w^\pm$ that are annihilated by $L_{m>0}, G_{n>0}$ and carries $L_0$ eigenvalue $h$. They are related by acting $G_0$ (we use the convention of \cite{0810.1203} for R-supermultiplet):
\ie
G_0 w^\pm = i \beta e^{\mp i\pi/4} w^\mp,
\fe
where $\beta$ is defined by $\beta^2 = \frac{c}{24}-h$. For R-sector $c$-recursion relation, instead of fixing the internal weight $h$, we fix $\beta$ of the internal primary, as doing so keeps the $G_0$ action on $w^\pm$ invariant. We take $w^+$ to be Grassmann even, and hence $w^-$ will be Grassmann odd. An arbitrary state in the Ramond multiplet can be then written as
\ie
\mathbb{L}_{-A} w^\pm \equiv L_{-n_1}\cdots L_{-n_k} G_{-m_1}\cdots G_{-m_\ell} w^\pm,\quad |A|= \sum_{i=1}^k n_i + \sum_{j=1}^\ell m_j,\quad (-1)^A \equiv (-1)^\ell,
\fe
where $n_1\geq \cdots \geq n_k, m_1>\cdots > m_\ell$. Throughout this notes, the NS-sector quantities will be written in \texttt{mathbf}, and the R-sector quantities in \texttt{mathbb}.

For a sphere with (NS, R, R) insertions, we normalize the 3-point function such that
\ie
\rho^{++}(\phi_1, w^+_2, w_3^+) = 1,\quad \rho^{+-}(\phi_1, w^+_2, w_3^-) = 1,\quad \rho^{-+}(\phi_1, w^-_2, w_3^+) = 1,\quad \rho^{--}(\phi_1, w^-_2, w_3^-) = 1.
\fe
The 3-point functions with other slot orders can be related to these basic ones by a Mobius transformation. It turns out that all $\rho^{\pm, \pm}$'s with all slot orders are normalized to 1. One can introduce \cite{1012.2974}
\ie
\rho_{0}^{(\eta)}(\xi_1, \xi_2, \xi_3) &= \rho^{++}(\xi_1, \xi_2, \xi_3) + \eta \rho^{--}(\xi_1, \xi_2, \xi_3),\\ \rho_{1}^{(\eta)}(\xi_1, \xi_2, \xi_3) &= \rho^{+-}(\xi_1, \xi_2, \xi_3) + i\eta \rho^{-+}(\xi_1, \xi_2, \xi_3),
\fe
for $\eta = \pm$. The subscript $0,1$ indicates the total fermion parity, as $\rho(\mathbf{L}_{-A}\phi, \mathbb{L}_{-C}w^\alpha, \mathbb{L}_{-E}w^\gamma)$ can eventually be reduced to $\rho(\phi, w^{\alpha'}, w^{\gamma'})$ with $(-1)^{|\alpha'|+|\gamma'|} = (-1)^{A+C+E+|\alpha'|+|\gamma'|}$, where $|\alpha|=0$ for $\alpha = +$ and $|\alpha|=1$ for $\alpha = -$, as the total fermion parity stays the same under the contour pulling used in Ward identities. 

The Ramond block in the glasses channel, with spin structure specified by the three signs $(\eta_1, \eta_2, \eta_3)$ and the edge 1 being NS, is defined by
\ie \label{Rblock}
\mathbb{F}_f^{(\eta, \eta')}(h_1, \beta_2, \beta_3, c) &= \sum_{\substack{|A|=|B|, |C|=|D|, |E|=|F|\\ (-1)^{A+C+E+|\alpha|+|\gamma|}=(-1)^f\\(-1)^{B+D+F+|\alpha'|+|\gamma'|}=(-1)^f}} q_1^{|A|} \eta_1^A q_2^{|C|} \eta_2^{C+|\alpha|} q_3^{|E|} \eta_3^{E+|\gamma|} \mathbf{G}^{AB}_{h_1} \mathbb{G}^{C\alpha, D\alpha'}_{\beta_2} \mathbb{G}^{E\gamma, F\gamma'}_{\beta_3}\\
&\quad\quad \times (-1)^{A(C+|\alpha|)+A(E+|\gamma|)+(C+|\alpha|)(E+|\gamma|)}\\
&\quad\quad \times \rho_f^{(\eta)}(\mathbf{L}_{-A}\phi_1, \mathbb{L}_{-C}w^\alpha_2, \mathbb{L}_{-E}w^\gamma_3)\rho_f^{(\eta')}(\mathbf{L}_{-B}\phi_1, \mathbb{L}_{-D}w^{\alpha'}_2, \mathbb{L}_{-F}w^{\gamma'}_3)
\
\fe
where $\mathbb{G}$ is the inverse Gram matrix for R-multiplets (which excludes actions of $G_0$), and $\alpha,\alpha', \gamma, \gamma'$ are such that the total fermion parity $(-1)^{A+C+E+|\alpha|+|\gamma|}=(-1)^f$. 

% One should note that since $\eta_i^F \mathbb{L}_{-A}w^+ = \eta_i^A L_{-A}w^+$ and $\eta_i^F \mathbb{L}_{-A}w^+ = -\eta_i^A L_{-A}w^+$, the action of $\eta_i^F$ will change the $\eta$-label of the 3-point function. We then have
% \ie
% \rho_f^{(\eta)}(\mathbf{L}_{-A} \phi_i, \eta_i^F \mathbb{L}_{-C} w^\alpha, \mathbb{L}_{-E} w^{\alpha^{\prime}})&=\eta_i^C \rho_f^{\left(\eta_i \eta\right)}(\mathbf{L}_{-A} \phi_i, \mathbb{L}_{-C} w^\alpha, \mathbb{L}_{-E} w^{\alpha^{\prime}}) ,\\
% \rho_f^{(\eta)}(\mathbf{L}_{-A} \phi_i, \mathbb{L}_{-C} w^\alpha, \eta_i^F \mathbb{L}_{-E} w^{\alpha^{\prime}})&=\eta_i^{E+f} \rho_f^{(\eta_i \eta)}(\mathbf{L}_{-A} \phi_i, \mathbb{L}_{-C} w^\alpha, \mathbb{L}_{-E} w^{\alpha^{\prime}}),
% \fe
% where the additional sign $\eta_i^f$ on the second line comes from the fact that $\rho_1^{(\eta)}=\rho^{+-}+i\eta\rho^{-+}$, hence the third-slot ground state is odd in the $+-$ component. 

Let's consider the case where the $w_2$ multiplet develops a null descendant. This happens when 
\ie
\beta=\beta_{r s}(c)=\frac{r b+s b^{-1}}{2 \sqrt{2}}, \quad r+s\in 2 \mathbb{Z}+1
\fe
and for fixed $\beta$, this becomes the following value for $c$:
\ie
x_{r s}^{\mathrm{R}}(\beta) & =\frac{4 \beta^2-r s+2 \sqrt{2} \beta \sqrt{2 \beta^2-r s}}{r^2}, \\
c_{r s}^{\mathrm{R}}(\beta) & =\frac{15}{2}+3\left(x_{r s}^{\mathrm{R}}(\beta)+x_{r s}^{\mathrm{R}}(\beta)^{-1}\right) .
\fe
We choose $b_{rs}=\sqrt{x_{rs}^{\mathrm{R}}}$ continuously such that $\beta=\beta_{rs}(c_{rs}^{\mathrm{R}})$. Below, $b_{rs}$ and $c_{rs}^{\mathrm{R}}$ refer to this branch. When this happens, the R multiplet develops a doublet of null states $\chi_{rs}^\pm = \sum_{M,\alpha} \chi_{rs}^{M,\alpha} G_0^\alpha \mathbb{L}_{-M} w^\pm$ happens at level $h+rs/2$, and we normalize $\chi_{rs}^\pm$ such that the coefficient of $L_{-1}^{rs/2}w^\pm$ is 1. In $\beta$ variables, the weight $h_{rs}=h+rs/2$ translates to
\ie
\beta_{r s}^{\prime}=\frac{(-1)^s}{2 \sqrt{2}}\left(r b_{r s}-s b_{r s}^{-1}\right).
\fe

The Ramond block $\mathbb{F}^{(\eta, \eta')}_f(h_1, \beta_2, \beta_3, c)$ develops a pole when $c=c_{rs}^{\mathrm{R}}(\beta_2)$. The inverse norm for the null state is given by
\ie
\mathbb{A}_{r s}^c \equiv \lim _{c \rightarrow c_{r a}^{\mathrm{R}}(\beta)} \frac{\beta_{r s}(c)^2-\beta^2}{\left\|\chi_{r s}^{+}\right\|^2}=2^{r s-2}\left(r b_{r s}+s b_{r s}^{-1}\right) \prod_{\substack{p=1-r, \ldots, r, q=1-s, \ldots, s \\ p+q \in \mathbb{Z},(p, q) \neq(0,0)}}\left(p b_{r s}+q b_{r s}^{-1}\right)^{-1}
\fe
and the fusion polynomials are given by 
\ie
\mathbb{P}_{rs}^{\R,\eta}(h_i,\beta_j;c)
={}&
\prod_{\substack{0\le k\le r-1,\ 0\le\ell\le s-1\\k+\ell\in2\mathbb Z}}
\frac{\lambda_i-2\sqrt2\,\eta\beta_j
 -(1-r+2k)b-(1-s+2\ell)b^{-1}}{2\sqrt2}\\
&\times
\prod_{\substack{0\le k\le r-1,\ 0\le\ell\le s-1\\k+\ell\in2\mathbb Z+1}}
\frac{\lambda_i+2\sqrt2\,\eta\beta_j
 -(1-r+2k)b-(1-s+2\ell)b^{-1}}{2\sqrt2}.
\fe
The factorization relations are summarized in appendix A. Using the relation
\ie
\rho_f^{(\eta)}
 (\mathbf{L}_{-A}\phi_i,
  \mathbb{L}_{-C}\chi_{rs}^{\alpha},
  \mathbb{L}_{-E}w_{\beta_j}^{\gamma})=
 \mathbb{P}_{rs}^{\R,\eta}(h_i,\beta_j;c_{rs}^{\R})
 \rho_f^{(\eta)}
 (\mathbf{L}_{-A}\phi_i,
  \mathbb{L}_{-C}w_{\beta'_{rs}}^{\alpha},
  \mathbb{L}_{-E}w_{\beta_j}^{\gamma})
\fe
one can read off the polar value of $\mathbb{F}_f^{(\eta, \eta')}$ at $c=c_{rs}^{\R}(\beta_2)$:
\ie
\operatorname{Res}_{c=c_{rs}^{\R}(\beta_2)} \mathbb{F}_f^{(\eta, \eta')} &= \frac{\partial c_{rs}(\beta_2)}{\partial \beta^2_2}\mathbb{A}_{rs}^c\sum_{\substack{|A|=|B|, |C|=|D|, |E|=|F|\\ (-1)^{A+C+E+|\alpha|+|\gamma|}=(-1)^f\\(-1)^{B+D+F+|\alpha'|+|\gamma'|}=(-1)^f}} q_1^{|A|}\eta_1^{A} q_2^{|C|}\eta_2^{C+|\alpha|} q_3^{|E|} \eta_3^{E+|\gamma|} \\
&\quad\quad \times (-1)^{A(C+|\alpha|)+A(E+|\gamma|)+(C+|\alpha|)(E+|\gamma|)}\mathbf{G}^{AB}_{h_1} \mathbb{G}^{C\alpha, D\alpha'}_{\beta'_{rs}} \mathbb{G}^{E\gamma, F\gamma'}_{\beta_3}\\
&\quad\quad \times \rho_f^{(\eta)}(\mathbf{L}_{-A}\phi_1, \mathbb{L}_{-C}\chi^\alpha_{rs}, \mathbb{L}_{-E}w^\gamma_3)\rho_f^{(\eta')}(\mathbf{L}_{-B}\phi_1, \mathbb{L}_{-D}\chi^\alpha_{rs}, \mathbb{L}_{-F}w^{\gamma'}_3),
\fe
and hence
\ie
\operatorname{Res}_{c=c_{rs}^{\R}(\beta_2)} \mathbb{F}_f^{(\eta, \eta')}&=\frac{\partial c_{rs}(\beta_2)}{\partial \beta^2_2}\mathbb{A}_{rs}^c\mathbb{P}_{rs}^{\R,\eta}(h_1,\beta_3;c_{rs}^\R(\beta_2))\mathbb{P}_{rs}^{\R,\eta'}(h_1,\beta_3;c_{rs}^\R(\beta_2))\\
&\quad \times\mathbb{F}_f^{(\eta, \eta')}(\beta_2\rightarrow \beta'_{rs},c\rightarrow c^{\R}_{rs}(\beta_2))
\fe
where we used the fact that $\chi_{rs}^{M,\alpha} G_0^\alpha \mathbb{L}_{-M}$ is Grassmann even, which ensures that the total fermion parity $f$ doesn't get shifted and there are no additional signs $\eta_2$. The residue at $c=c^{\R}_{rs}(\beta_3)$ can be worked out analogously.

The Ramond also develops poles when the NS supermultiplet contains a null state. The relavant fusion polynomials are
\ie
\mathbb{P}_{rs}^{\NS,\eta}(\beta_2,\beta_3;c)
={}&
\prod_{\substack{0\le k\le r-1,\ 0\le\ell\le s-1\\k+\ell\in2\mathbb Z}}
\frac{2\sqrt2(\beta_3-\eta\beta_2)
 -(1-r+2k)b-(1-s+2\ell)b^{-1}}{2\sqrt2}\\
&\times
\prod_{\substack{0\le k\le r-1,\ 0\le\ell\le s-1\\k+\ell\in2\mathbb Z+1}}
\frac{2\sqrt2(\beta_3+\eta\beta_2)
 -(1-r+2k)b-(1-s+2\ell)b^{-1}}{2\sqrt2},
\label{eq:appendix-ns-rr-fusion-polynomial}
\fe
and the factorization relation is given by
\ie
&\rho_f^{(\eta)}
 (\mathbf{L}_{-A}\chi_{rs},
  \mathbb{L}_{-C}w_{\beta_2}^{\alpha},
  \mathbb{L}_{-E}w_{\beta_3}^{\gamma})\\
&\quad=\mathbb{P}_{rs}^{\NS,\eta}(\beta_2,\beta_3;c_{rs})
\begin{cases}
\rho_f^{(\eta)}
 (\mathbf{L}_{-A}\phi_{h_{rs}+rs/2},
  \mathbb{L}_{-C}w_{\beta_2}^{\alpha},
  \mathbb{L}_{-E}w_{\beta_3}^{\gamma}),&rs\in2\mathbb Z,\\[0.4em]
e^{i(-1)^f\pi/4}\rho_{1-f}^{(-\eta)}
 (\mathbf{L}_{-A}\phi_{h_{rs}+rs/2},
  \mathbb{L}_{-C}w_{\beta_2}^{\alpha},
  \mathbb{L}_{-E}w_{\beta_3}^{\gamma}),&rs\in2\mathbb Z+1.
\end{cases}
\fe
From these, one can derive the residue at $c=c_{rs}(h_1)$:
\ie
&\operatorname{Res}_{c=c_{r s}\left(h_1\right)} \mathbb{F}_f^{\left(\eta, \eta^{\prime}\right)} = -\frac{\partial c_{r s}\left(h_1\right)}{\partial h_1} \mathbf{A}_{r s}^c(h_1) q_1^{r s / 2} \mathbb{P}_{r s}^{\mathrm{NS}, \eta}\left(\beta_2, \beta_3 ; c_{r s}(h_1)\right) \mathbb{P}_{r s}^{\mathrm{NS}, \eta^{\prime}}\left(\beta_2, \beta_3 ; c_{r s}(h_1)\right) \\
&\qquad \times \begin{cases}
\mathbb{F}_f^{\left(\eta, \eta^{\prime}\right)}\left(h_1\rightarrow h_1+ \tfrac{rs}{2}, c\rightarrow c_{r s}(h_1)\right),  & rs\in 2\mathbb{Z}, \\
i(-1)^f \eta_1 \mathbb{F}_{1-f}^{\left(-\eta, -\eta^{\prime}\right)}\left(h_1\rightarrow h_1+ \tfrac{rs}{2}, c\rightarrow c_{r s}(h_1),\eta_2\rightarrow -\eta_2, \eta_3\rightarrow -\eta_3\right), & rs\in 2\mathbb{Z}+1.
\end{cases}
\fe

We now have computed all the polar part in the following recursion relation
\ie
\mathbb{F}^{(\eta, \eta')}_f(h_1, \beta_2, \beta_3, c) &= \mathbb{F}^{(\eta, \eta')}_f(h_1, \beta_2, \beta_3, c\rightarrow \infty)\\
&+ \sum_{r+s\in 2\mathbb{Z}} \frac{\operatorname{Res}_{c=c_{rs}(h_1)}\mathbb{F}^{(\eta, \eta')}_f}{c-c_{rs}(h_1)} + \sum_{i=2,3}\sum_{r+s\in 2\mathbb{Z}+1} \frac{\operatorname{Res}_{c=c_{rs}^{\R}(\beta_i)} \mathbb{F}_f^{(\eta, \eta')}}{c-c_{rs}^{\R}(\beta_i)}
\fe
The next ingredient of the $c$-recursion relation of the Ramond blocks is the large-$c$ limit, which we will compute in the next section.

\section{The large-$c$ Ramond block}

In the $c\rightarrow \infty$ limit, we fix $h$ of the NS supermultiplet, and fix $\beta$ of the R supermultiplet, which implies that $L_0 = \frac{c}{24} - \beta^2 + |A|$ is of order $\mathcal{O}(c)$ for a Ramond state $\mathbb{L}_{-A}w^\pm_\beta$.

We first analyze the $c$-scaling behavior of the Gram matrix element $\langle w^\alpha_\beta|\mathbb{L}_{-A}^\dagger \mathbb{L}_{-B}|w^\gamma_\beta\rangle$, where there are no $G_0$'s in the lowering operators $\mathbb{L}_{-A}$. The SCA in the R-sector reads
\ie
{\left[L_m, L_n\right] } & =(m-n) L_{m+n}+\frac{c}{12}\left(m^3-m\right) \delta_{m+n, 0}, \\
{\left[L_m, G_n\right] } & =\left(\frac{m}{2}-n\right) G_{m+n}, \\
\left\{G_m, G_n\right\} & =2 L_{m+n}+\frac{c}{3}\left(m^2-\frac{1}{4}\right) \delta_{m+n, 0}.
\fe
Therefore, in the large-$c$ limit, the algebra becomes
\ie
{\left[L_m, L_{-n}\right] } =\frac{c}{12}m^3 \delta_{m, n},\quad \left\{G_m, G_{-n}\right\} = \frac{c}{3}m^2 \delta_{m, n},\quad {\left[L_m, G_{-n}\right] } = 0.
\fe
The leading term in the Gram matrix then comes from pairwise contractions between the generators. Let $\hat{A}$ be the total number of generators in $\mathbb{L}_{-A}$, then $\langle w^\alpha_\beta|\mathbb{L}_{-A}^\dagger \mathbb{L}_{-B}|w^\gamma_\beta\rangle$ is of order $\mathcal{O}(c^{\hat{A}})$ if $(A,\alpha)=(B,\gamma)$, at most $\mathcal{O}(c^{\hat{A}-1})$ if $\hat{A}=\hat{B}, (A,\alpha)\neq (B,\gamma)$, and at most $\mathcal{O}(c^{\mathrm{min}(\hat{A}, \hat{B})})$ if $\hat{A}\neq \hat{B}$. One can then apply Gram-Schmidt to $\mathbb{L}_{-A}w^{\pm}_\beta$ to obtain a orthonormal basis $\mathcal{L}_{-A}w^{\pm}_\beta$, and
\ie
\|\mathcal{L}_{-A}w^\pm_{\beta}\|^2 \sim \mathcal{O}(c^{\hat{A}}).
\fe

The next ingredient is the scaling behavior of the 3-point function $\rho_f^{(\eta)}(\mathbf{L}_{-A}\phi_1, \mathbb{L}_{-B}w^\alpha_2,\mathbb{L}_{-C}w^\gamma_3)$. One can use conformal Ward identities to convert a $L_{-n}, G_{-r}$ insertion in one slot to another slot. For a $L_{-n}$ insertion,
\ie
& \rho_f^{(\eta)}\left(L_{-n} \mathbf{L}_{-A} \phi_1, \mathbb{L}_{-B} w_2^\alpha, \mathbb{L}_{-C} w_3^\gamma\right) \\
& \qquad=\rho_f^{(\eta)}\left(\mathbf{L}_{-A} \phi_1, \mathbb{L}_{-B} w_2^\alpha, L_n \mathbb{L}_{-C} w_3^\gamma\right)+\sum_{p=-1}^n\binom{n+1}{p+1} \rho_f^{(\eta)}\left(\mathbf{L}_{-A} \phi_1, L_p \mathbb{L}_{-B} w_2^\alpha, \mathbb{L}_{-C} w_3^\gamma\right),\\
& \rho_f^{(\eta)}\left(\mathbf{L}_{-A} \phi_1, \mathbb{L}_{-B} w_2^\alpha, L_{-n} \mathbb{L}_{-C} w_3^\gamma\right) \\
& \qquad=\rho_f^{(\eta)}\left(L_n \mathbf{L}_{-A} \phi_1, \mathbb{L}_{-B} w_2^\alpha, \mathbb{L}_{-C} w_3^\gamma\right)-\sum_{p=0}^{\infty}\binom{1-n}{p} \rho_f^{(\eta)}\left(\mathbf{L}_{-A} \phi_1, L_{p-1} \mathbb{L}_{-B} w_2^\alpha, \mathbb{L}_{-C} w_3^\gamma\right) ,\\
&\rho_f^{(\eta)}  \left(\mathbf{L}_{-A} \phi_1, L_{-n} \mathbb{L}_{-B} w_2^\alpha, \mathbb{L}_{-C} w_3^\gamma\right) = \sum_{p=0}^{\infty}\binom{n-2+p}{n-2} \rho_f^{(\eta)}\left(L_{n+p} \mathbf{L}_{-A} \phi_1, \mathbb{L}_{-B} w_2^\alpha, \mathbb{L}_{-C} w_3^\gamma\right) \\
& \qquad+(-1)^n \sum_{p=0}^{\infty}\binom{n-2+p}{n-2} \rho_f^{(\eta)}\left(\mathbf{L}_{-A} \phi_1, \mathbb{L}_{-B} w_2^\alpha, L_{p-1} \mathbb{L}_{-C} w_3^\gamma\right) .
\fe
This imposes a problem to our infinite-$c$ limit. Each conformal Ward identity contains a $L_0$ term acting on a Ramond multiplet on the RHS. In our fixed $\beta$ large-$c$ limit, this will contribute to a term of order $\mathcal{O}(c)$. Therefore, for $\mathbf{L}_{-A}, \mathbb{L}_{-A}$ made from purely $L_{-n}$'s,
\ie
\rho_f^{(\eta)}(\mathbf{L}_{-A}\phi_1, \mathbb{L}_{-B}w^\alpha_2,\mathbb{L}_{-C}w^\gamma_3) \sim \mathcal{O}(c^{\hat{A}+\hat{B}+\hat{C}})
\fe
and the contribution of this term to the large-$c$ conformal block will scale like $\mathcal{O}(c^{\hat{A}+\hat{B}+\hat{C}})$. Namely, the large-$c$, fixed-$\beta$ limit of the Ramond block is not well-defined. 

\section{NS blocks from double Virasoro}

We now turn to another approach of computing the superconformal blocks. The idea is: Given the $\mathcal{N}=1$ SCA with shift $\nu\in \left\{0,\frac{1}{2}\right\}$ in the fermionic level, one can construct an enlarged algebra $\mathsf{SCA}_\nu \otimes \mathsf{F}_{\nu}$, where $\mathsf{F}_{\nu}$ a auxiliary free fermion algebra with the same level shift $\nu$.  $\mathsf{SCA}_\nu \otimes \mathsf{F}_{\nu}$ admits two copy of virasoro algebra as it's subalgebra $\mathsf{Vir}\oplus \mathsf{Vir}\subset \mathsf{SCA}_\nu \otimes \mathsf{F}_{\nu}$ \cite{1312.4520,1608.02568}. Using the branching rules of this subalgebra reduction, one can decompose the superconformal block into ordinary Virasoro blocks.

The enlarged algebra $\mathsf{SCA}_\nu \otimes \mathsf{F}_{\nu}$ is given by:
\ie
{\left[L_m, L_n\right] } & =(m-n) L_{m+n}+\frac{c}{12}\left(m^3-m\right) \delta_{m+n, 0}, \\
{\left[L_m, G_r\right] } & =\left(\frac{m}{2}-r\right) G_{m+r}, \\
\left\{G_r, G_s\right\} & =2 L_{r+s}+\frac{c}{3}\left(r^2-\frac{1}{4}\right) \delta_{r+s, 0} ,\\
\{\psi_r, \psi_s\} &= \delta_{r+s,0},\quad \{\psi_r, G_s\}=[\psi_r, L_n] = 0.
\fe
where the levels $m,n\in\mathbb{Z}$, $r, s\in \mathbb{Z}+\nu$. The reality conditions are
\ie
\psi_{r}^\dagger = -\psi_{-r}, \quad G_{r}^\dagger = G_{-r}, \quad L_{n}^\dagger = L_{-n}.
\fe
We also introduce the following operators
\ie
L_n^{\mathrm{F}} = \frac{1}{2} \sum_{r\in \mathbb{Z}+\nu} r:\psi_{n-r}\psi_r: + \frac{1-2\nu}{16}\delta_{n,0},\quad U_n = \sum_{r\in \mathbb{Z}+\nu} \psi_{n-r} G_r.
\fe
Define $b$ by the relation
\ie
c= \frac{3}{2}+3\left(b+b^{-1}\right)^2.
\fe
For $b\neq 1$, one can then define
\ie
\begin{aligned}
& L_n^{(1)}=\frac{b^{-1}}{b^{-1}-b} L_n-\frac{b^{-1}+2 b}{b^{-1}-b} L_n^{\mathrm{F}}+\frac{1}{b^{-1}-b} U_n, \\
& L_n^{(2)}=\frac{b}{b-b^{-1}} L_n-\frac{b+2 b^{-1}}{b-b^{-1}} L_n^{\mathrm{F}}+\frac{1}{b-b^{-1}} U_n ,
\end{aligned}
\fe
which satisfies $L_n^{(1)}+L_n^{(2)}=L_n+L_n^{\mathrm{F}}$. Note that in order to have $(L_n^{(i)})^\dagger = L_{-n}^{(i)}$, we need to take the BPZ conjugate of $\psi_r$ to be $(\psi_r)^\dagger = -\psi_{-r}$. One can also prove that these operators are the generators of two commuting copies of the Virasoro algebra
\ie
[L_m^{(i)}, L_n^{(j)}] = \left[(m-n) L_{m+n}^{(i)} + \frac{c^{(i)}}{12} (m^3-m)\delta_{m+n,0}\right]\delta_{ij},
\fe
and the two central charges are given by
\ie
c^{(i)}=1+6\left(b^{(i)}+\frac{1}{b^{(i)}}\right)^2,\quad (b^{(1)})^2 = \frac{2b^2}{1-b^2},\quad (b^{(2)})^{-2} = \frac{2}{b^2-1}.
\fe
This proves the subalgebra relation $\mathsf{Vir}\oplus \mathsf{Vir}\subset \mathsf{SCA}_\nu \otimes \mathsf{F}_{\nu}$. 

Under the subalgebra reduction $\mathsf{Vir}\oplus \mathsf{Vir}\subset \mathsf{SCA}_\nu \otimes \mathsf{F}_{\nu}$, a highest-weight representation of $\mathsf{SCA}_\nu \otimes \mathsf{F}_{\nu}$ can be decomposed into irreps of $\mathsf{Vir}\oplus \mathsf{Vir}$. This decomposition is worked out in \cite{1608.02568}. We now discuss this decomposition in both NS and R sectors.

\subsubsection*{NS Sector}
As always, let $\mathcal{V}_{\mathsf{SCA}_{1/2}}(h)$ denote the NS Verma
module generated by an  highest-weight state $\phi$ of conformal weight $h$. Let $\mathcal{V}_{F_{1/2}}$ denote the Fock representation of $F_{1/2}$. The representation of interest is $\mathcal{M}_{\mathrm{NS}}(h):=\mathcal{V}_{F_{1/2}}\otimes \mathcal{V}_{\mathsf{SCA}_{1/2}}(h)$.
It can be generated by acting arbitrary $\psi_{-r}, L_{-n}, G_{-r}$ on $\mathbf{1}_{\mathsf{F}}\otimes \phi$. 

We further define the "Liouville momentum" $P$ by
\ie
h=\frac{1}{2}\left(\frac{1}{4}Q^2-P^2\right),
\fe
where $Q=b+b^{-1}$.  Denoting the Virasoro Verma module with lowest weight $h$ by $\mathcal{V}_c(h)$, it's proven that \cite{1111.2803}
\ie \label{NSbranching}
\mathcal{M}_{\mathrm{NS}}(h) \cong \bigoplus_{n \in \frac{1}{2} \mathbb{Z}} \mathcal{V}_{c^{(1)}}(h_{n}^{(1)}) \otimes \mathcal{V}_{c^{(2)}}(h_{n}^{(2)})
\fe
for $P\notin \frac{1}{2}(b\mathbb{Z}+b^{-1}\mathbb{Z})$. The weights $h^{(i)}_{n}$'s are given by $h^{(i)}_n = \frac{1}{4}(Q^{(i)})^2-(P_{n}^{(i)})^2$, where
\ie \label{branchweights}
P_n^{(1)}=\frac{P}{\sqrt{2-2 b^2}}+n b^{(1)},\quad P_n^{(2)}=\frac{P}{\sqrt{2-2 b^{-2}}}+n(b^{(2)})^{-1}.
\fe
It satisfies $h_{n}^{(1)}+h_{n}^{(2)}= h+2n^2$. 

We will denote the highest weight vector in $\mathcal{V}_{c^{(1)}}(h_{n}^{(1)}) \otimes \mathcal{V}_{c^{(2)}}(h_{n}^{(2)})$ by $v_{n}$. In our convention where $\phi$ is Grassmann even, $v_n$ has Grassmann parity $2n$. The Liouville momentum $P$ gives identical $\mathcal{V}_{\mathsf{SCA}_{1/2}}$, and a reflection $(n, P)\mapsto (-n, -P)$ gives the same $\mathsf{Vir}\oplus \mathsf{Vir}$ module. We can then identify $v_n(P)=v_{-n}(-P)$ and work primarily with $P>0$. For multiplets with $h > Q^2/8$, we analytic continue all the formula from positive real $P$ to purely imaginary $P=ip$ with $p>0$. 

For our application, we will always use the following bilinear pairing for $\mathcal{M}_{\mathrm{NS}}(h)$
\ie
\bpzbraket{\mathbf{\Psi}_{-\mathsf{A}}\mathbf{1}_F \otimes \mathbf{L}_{-A} \phi }{\mathbf{\Psi}_{-\mathsf{B}}\mathbf{1}_F \otimes \mathbf{L}_{-B} \phi}^{\prime} =\bpzbraket{\mathbf{\Psi}_{-\mathsf{A}}\mathbf{1}_F}{\mathsf{\Psi}_{-\mathbf{B}}\mathbf{1}_F} \times \bpzbraket{\mathbf{L}_{-A} \phi }{\mathbf{L}_{-B} \phi}^{\prime}.
\fe
where we use $A+B = 0~\text{mod}2$ for the matrix element to be non-vanishing. This bilinear pairing has the property that 
\ie
\bpzbraket{\mathbf{L}^{(1)}_{-M} \nu_1\otimes \mathbf{L}^{(2)}_{-N}  \nu_2} { \mathbf{L}^{(1)}_{-M'} \otimes \nu_1\mathbf{L}^{(2)}_{-N'} \otimes \nu_2}^{\prime} = \bpzbraket{\mathbf{L}^{(1)}_{-M}  \nu_1} {\mathbf{L}^{(1)}_{-M'} \nu_1}\bpzbraket{\mathbf{L}^{(2)}_{-N}  \nu_2} {\mathbf{L}^{(2)}_{-N'} \nu_2}^{\prime}.
\fe
It should be distinguished from the standard bilinear pairing on $\mathcal{M}_{\mathrm{NS}}(h)$ in the theory of chiral fermion coupled to superconformal field theory, where one need an additional crossing sign
\ie
\bpzbraket{\mathbf{\Psi}_{-\mathsf{A}}\mathbf{1}_F \otimes \mathbf{L}_{-A} \phi }{\mathbf{\Psi}_{-\mathsf{B}}\mathbf{1}_F \otimes \mathbf{L}_{-B} \phi} =(-1)^{\mathsf{B}(A+p_{\phi} )}\bpzbraket{\mathbf{\Psi}_{-\mathsf{A}}\mathbf{1}_F}{\mathsf{\Psi}_{-\mathbf{B}}\mathbf{1}_F} \times \bpzbraket{\mathbf{L}_{-A} \phi }{\mathbf{L}_{-B} \phi}.
\fe


\subsubsection*{R Sector}
Let $\mathcal{V}_{\mathsf{SCA}_{0}}(\beta)$ denote the R Verma
module generated by an highest-weight state $w^{+}_{\beta}$ of conformal weight $h = \frac{c}{24}-\beta^2$. Two degenerate ground states will be denoted as $w^{\pm}_{\beta}$. Let $\mathcal{V}_{F_0}$ denote the Fock representation of $F_{0}$. Two ground states in $\mathcal{V}_{F_0}$ will be denoted as $u^\pm$, with $\psi_0 u^\pm = 2^{-1/2}u^\mp$.

The representation of interest is $\mathcal{M}_{\mathrm{R}}(\beta):= \mathcal{V}_{F_{0}}\otimes  \mathcal{V}_{\mathsf{SCA}_{0}}(\beta)$. It contains 4 ground states given by $u^\pm \otimes w^\pm_\beta$. $\mathcal{M}_{\mathrm{R}}(\beta)$ can be generated by acting lowering operators on the 4 ground states $u^\pm \otimes w_\beta^\pm$. The Liouville momentum of $w_\beta^\pm$ is defined by
\ie
h = \frac{c}{24}-\beta^2 = \frac{1}{16} + \frac{1}{2}\left(\frac{1}{4}Q^2-P^2\right),\quad P^2 = 2\beta^2.
\fe
For $P\notin \frac{1}{2}(b\mathbb{Z}+b^{-1}\mathbb{Z})$, $\mathcal{M}_{\mathrm{R}}(\beta)$ can be decomposed into \cite{1608.02568}
\ie\label{Rbranching}
\mathcal{M}_{\mathrm{R}}(\beta) \cong \bigoplus_{\epsilon=0,1} \bigoplus_{n \in \frac{1}{2} \mathbb{Z}+\frac{1}{4}} \mathcal{V}_{c^{(1)}}(h_{n}^{(1)}) \otimes \mathcal{V}_{c^{(2)}}(h_{n}^{(2)}),
\fe
where
\ie
h^{(i)}_n = \frac{1}{4}(Q^{(i)})^2-(P_{n}^{(i)})^2,\quad P_n^{(1)}=\frac{P}{\sqrt{2-2 b^2}}+n b^{(1)},\quad P_n^{(2)}=\frac{P}{\sqrt{2-2 b^{-2}}}+n(b^{(2)})^{-1}.
\fe
It satisfies $h_n^{(1)}+h_n^{(2)} = h +2n^2 -\frac{1}{16}$.

The direct sum over $\epsilon$ indicates that there are two copies of each $\mathcal{V}_{c^{(1)}}(h_{n}^{(1)}) \otimes \mathcal{V}_{c^{(2)}}(h_{n}^{(2)})$, and $\epsilon$ denotes the total fermion parity of the lowest weight state (including the $\psi$ parity). One can directly check that the characters of the two sides of \eqref{NSbranching}, \eqref{Rbranching} agrees via the Jacobi triple product identity. The ground state of $\mathcal{V}_{c^{(1)}}(h_{n}^{(1)}) \otimes \mathcal{V}_{c^{(2)}}(h_{n}^{(2)})$ will be denoted by $v^\pm_n$. We use the convention that $v^\alpha_n$ has Grassmann parity $\alpha$.

\subsubsection*{Decomposition of all-NS block}

Let's inspect on how the superconformal blocks $\mathbf{F}_a(h_1, h_2, h_3, c)$, $\mathbb{F}_f^{(\eta, \eta')}(h_1, \beta_2, \beta_3,c)$ decompose under $\mathsf{Vir}\oplus \mathsf{Vir}\subset \mathsf{SCA}_\nu \otimes \mathsf{F}$. We first define the all-NS $\mathsf{SCA}_\nu \otimes \mathsf{F}$ block in the theta channel as
\ie
\widehat{\mathbf{F}}_a(h_1, h_2, h_3, c)= & \sum_{
  \substack{
  |A|=|A'|,|B|=|B'|,|C|=|C'|,\mathsf{A}, \mathsf{B}, \mathsf{C}\\
  (-1)^{A+B+C}(-1)^{\mathsf{A}+\mathsf{B}+\mathsf{C}}=(-1)^a\\
  }
} q_1^{|A|}\eta_1^A q_2^{|B|} \eta_2^B  q_3^{|C|} \eta_3^C \mathbf{G}_{h_1}^{AA'} \mathbf{G}_{h_2}^{BB'} \mathbf{G}_{h_3}^{CC'} \\
& \times (-1)^{(A+\mathsf{A})(B+\mathsf{B})+(A+\mathsf{A})(C+\mathsf{C})+(B+\mathsf{B})(C+\mathsf{C})}(-1)^{\mathsf{A}+\mathsf{B}+\mathsf{C}}q_1^{|\mathsf{A}|}\eta_1^{\mathsf{A}}q_2^{|\mathsf{B}|}\eta_2^{\mathsf{B}}q_3^{|\mathsf{C}|} \eta_3^{\mathsf{C}}\\
& \times \hat{\rho}_a(\Psi_{-\mathsf{A}}\mathbf{L}_{-A} \phi_1, \Psi_{-\mathsf{B}}\mathbf{L}_{-B} \phi_2, \Psi_{-\mathsf{C}}\mathbf{L}_{-C} \phi_3) \hat{\rho}_a(\Psi_{-\mathsf{A}}\mathbf{L}_{-A'} \phi_1, \Psi_{-\mathsf{B}}\mathbf{L}_{-B'} \phi_2, \Psi_{-\mathsf{C}}\mathbf{L}_{-C'} \phi_3) ,
\fe
where ${\Psi}_{-\mathsf{A}}v$ stands for $\Psi_{-\mathsf{A}}v \equiv \psi_{-r_1}\cdots \psi_{-r_m} v$, $|\mathsf{A}|$ stands for the total level of fermion lowering operators, and $(-1)^{\mathsf{A}}$ is the total $\psi$ parity. We have used the fact that $\bpzbraket{\mathbf{L}_{-A}v}{\Psi_{-\mathsf{A}}^\dagger \Psi_{-\mathsf{B}}\big | \mathbf{L}_{-B}v}$ is the matrix $(-1)^{\mathsf{A}}\delta_{\mathsf{A},\mathsf{B}}\mathbf{G}_{AB}$, which give rises to the sign $(-1)^{\mathsf{A}+\mathsf{B}+\mathsf{C}}$. The 3-point function $\hat{\rho}^{L/R}$ is labeled by the total parity, including the fermionic one.


For states $u\otimes x$, the 3-point function can be factorized into
\ie
& \hat{\rho}_a\left(u_1 \otimes x_1, u_2 \otimes x_2, u_3 \otimes x_3\right)=(-1)^{\left|x_2\right|\left|u_3\right|} \rho^{\mathsf{F}}(u_1, u_2, u_3) \rho_a(x_1, x_2, x_3),
\fe
where $\rho^{\mathsf{F}}(u_1, u_2, u_3)$ is the normalized 3-point function in the free fermion theory.  The sign $(-1)^{\left|x_2\right|\left|u_3\right|}$ comes from moving operator order. Importantly, $\rho^{\mathsf{F}}$ is only nonzero if $\mathsf{A}+\mathsf{B}+\mathsf{C}\equiv 0$ mod 2. If we define the free fermion block as
\ie
\mathbf{F}^{\mathsf{F}}(q_i, \eta_i) =& \sum_{\substack{\mathsf{A}, \mathsf{B}, \mathsf{C}\\(-1)^{\mathsf{A}+\mathsf{B}+\mathsf{C}}=1}}(-1)^{\mathsf{AB}+\mathsf{AC}+\mathsf{BC}}q_1^{|\mathsf{A}|}\eta_1^{\mathsf{A}}q_2^{|\mathsf{B}|}\eta_2^{\mathsf{B}}q_3^{|\mathsf{C}|} \eta_3^{\mathsf{C}}\left[{\rho}^{\mathsf{F}}(\Psi_{-\mathsf{A}}\mathbf{1}_\mathsf{F}, \Psi_{-\mathsf{B}}\mathbf{1}_\mathsf{F}, \Psi_{-\mathsf{C}}\mathbf{1}_\mathsf{F})\right]^2,
\fe
then the all-NS $\mathsf{SCA}_\nu \otimes \mathsf{F}$ block can be factorized into
\ie
\widehat{\mathbf{F}}_a(h_1, h_2, h_3, c;q_i, \eta_i)=\mathbf{F}_{\mathsf{F}}(q_i, \eta_i)\star {\mathbf{F}}_a(h_1, h_2, h_3, c;q_i, \eta_i),
\fe
where the symbol $\star$ stands for convolution: Denote the coefficient of $q_1^{n_1/2}q_2^{n_2/2}q_3^{n_3/2}$ in the $q$-series $\mathbf{F}_a$ by $\mathbf{F}_a[q_1^{n_1/2}q_2^{n_2/2}q_3^{n_3/2}]$, this equation can be written as
\ie
\widehat{\mathbf{F}}_a(h_1, h_2, h_3, c)&[q_1^{n_1/2}q_2^{n_2/2}q_3^{n_3/2}]\\
&=\sum_{r_i+s_i=n_i}(-1)^{\sum_{i<j} (r_is_j+s_ir_j)}\mathbf{F}_{\mathsf{F}}[q_1^{r_1/2}q_2^{r_2/2}q_3^{r_3/2}] {\mathbf{F}}_a(h_1, h_2, h_3, c)[q_1^{s_1/2}q_2^{s_2/2}q_3^{s_3/2}]
\fe

Using the decomposition of the representation $\mathcal{M}_{\text{NS}}$ under $\mathsf{Vir}\oplus \mathsf{Vir}\subset \mathsf{SCA}_\nu \otimes \mathsf{F}$, we can also express $\widehat{\mathbf{F}}_a(h_1, h_2, h_3, c)$ in terms of Virasoro blocks. Using the fact that a basis of state of a $\mathsf{SCA}_{1/2}\otimes \mathsf{F}$ module is given by $\{L_{-A}^{(1)}L_{-B}^{(2)}v_n\}$, one have
\ie
&\widehat{\mathbf{F}}_a(h_1, h_2, h_3, c) = \sum_{\substack{|A|=|A'|, |B|=|B'|, |C|=|C'| \\ |D|=|D'|, |E|=|E'|, |F|=|F'| \\ n_1, n_2, n_3\in \frac{1}{2}\mathbb{Z} \\ (-1)^{2(n_1+n_2+n_3)}=(-1)^a}} q_1^{2n_1^2+|A|+|B|}q_2^{2n_2^2+|C|+|D|}q_3^{2n_3^2+|E|+|F|}\eta_1^{2n_1}\eta_2^{2n_2}\eta_3^{2n_3}\\
&\quad \times (-1)^{4n_1n_2+4n_1n_3+4n_2n_3} \frac{(G_{h_{1,n_1}^{(1)}})^{AA'}(G_{h_{1,n_1}^{(2)}})^{BB'}(G_{h_{2,n_2}^{(1)}})^{CC'}(G_{h_{2,n_2}^{(2)}})^{DD'}(G_{h_{3,n_3}^{(1)}})^{EE'}(G_{h_{3,n_3}^{(2)}})^{FF'}}{\|v_{1,n_1}\|^2\|v_{2,n_2}\|^2\|v_{3,n_3}\|^2}\\
&\quad \times \hat{\rho}_a(L_{-A}^{(1)}L_{-B}^{(2)}v_{1,n_1},L_{-C}^{(1)}L_{-D}^{(2)}v_{2,n_2},L_{-E}^{(1)}L_{-F}^{(2)}v_{3,n_3})\hat{\rho}_a(L_{-A'}^{(1)}L_{-B'}^{(2)}v_{1,n_1},L_{-C'}^{(1)}L_{-D'}^{(2)}v_{2,n_2},L_{-E'}^{(1)}L_{-F'}^{(2)}v_{3,n_3}),
\fe
where we have used the fact that $h_{n}^{(1)}+h_{n}^{(2)}-h = 2n^2$ and that $v_n$ has Grassmann parity $2n$. $G_h^{AB}$ is the Virasoro inverse Gram matrix. The resolution of $(-1)^{p_{\mathsf{SCA}}\cdot p_{\mathsf{F}}}$ leads to the factorialization of the Gram matrix.

The 3-point function $\hat{\rho}_a(L_{-A}^{(1)}L_{-B}^{(2)}v_{1,n_1},L_{-C}^{(1)}L_{-D}^{(2)}v_{2,n_2},L_{-E}^{(1)}L_{-F}^{(2)}v_{3,n_3})$ can be reduced to those of Virasoro primaries $\hat{\rho}_a(v_{1,n_1}, v_{2,n_2}, v_{3,n_3})$ by conformal Ward identity. The two copies of Virasoro can be reduced independently, which leads to the factorization formula
\ie
\hat{\rho}_a&(L_{-A}^{(1)}L_{-B}^{(2)}v_{1,n_1},L_{-C}^{(1)}L_{-D}^{(2)}v_{2,n_2},L_{-E}^{(1)}L_{-F}^{(2)}v_{3,n_3})\\
&\quad =\hat{\rho}_a(v_{1,n_1}, v_{2,n_2}, v_{3,n_3})\rho(L_{-A}\nu_{h_{1,n_1}^{(1)}}, L_{-C}\nu_{h_{2,n_2}^{(1)}}, L_{-E}\nu_{h_{3,n_3}^{(1)}})\rho(L_{-B}\nu_{h_{1,n_1}^{(2)}}, L_{-D}\nu_{h_{2,n_2}^{(2)}}, L_{-F}\nu_{h_{3,n_3}^{(2)}}),
\fe
where $\rho(L_{-A}\nu_1, L_{-B}\nu_2, L_{-C}\nu_3)$ is defined in section 1. If we define the branching coefficient
\ie
B_a(h_1, h_2, h_3;n_1, n_2, n_3) = \frac{\hat{\rho}_a(v_{1,n_1}, v_{2,n_2}, v_{3,n_3})}{\|v_{1,n_1}\|\|v_{2,n_2}\|\|v_{3,n_3}\|},
\fe
the $\mathsf{SCA}_\nu \otimes \mathsf{F}$ block can be decomposed into
\ie
&\widehat{\mathbf{F}}_a(h_1, h_2, h_3, c) = \sum_{\substack{n_1, n_2, n_3\in \frac{1}{2}\mathbb{Z} \\ (-1)^{2(n_1+n_2+n_3)}=(-1)^a}} q_1^{2n_1^2}q_2^{2n_2^2}q_3^{2n_3^2}\eta_1^{2n_1}\eta_2^{2n_2}\eta_3^{2n_3}(-1)^{4n_1n_2+4n_1n_3+4n_2n_3}\\
&\qquad \times [B_a(h_1,h_2,h_3;n_1,n_2,n_3)]^2 F(h_{1,n_1}^{(1)},h_{2,n_2}^{(1)},h_{3,n_3}^{(1)},c^{(1)})F(h_{1,n_1}^{(2)},h_{2,n_2}^{(2)},h_{3,n_3}^{(2)},c^{(2)}),
\fe
where $F(h_1, h_2, h_3, c)$ is the Virasoro conformal block, which can be computed efficiently using the $c$-relation relation in section 1. Since the fermion block $\mathbf{F}_{\mathsf{F}}$ will always have the identity term $1$, one can find $\mathbf{F}_{\mathsf{F}}^{-1}$ as a formal power series in $q$, such that $\mathbf{F}_{\mathsf{F}}\star \mathbf{F}_{\mathsf{F}}^{-1}=1$. The all-NS superconformal block $\mathbf{F}_a$ can then be computed as
\ie
&\mathbf{F}_a(h_1, h_2, h_3, c) = (\mathbf{F}_{\mathsf{F}})^{-1}\star \sum_{\substack{n_1, n_2, n_3\in \frac{1}{2}\mathbb{Z} \\ (-1)^{2(n_1+n_2+n_3)}=(-1)^a}} q_1^{2n_1^2}q_2^{2n_2^2}q_3^{2n_3^2}\eta_1^{2n_1}\eta_2^{2n_2}\eta_3^{2n_3}(-1)^{4n_1n_2+4n_1n_3+4n_2n_3}\\
&\qquad \times [B_a(h_1,h_2,h_3;n_1,n_2,n_3)]^2 F(h_{1,n_1}^{(1)},h_{2,n_2}^{(1)},h_{3,n_3}^{(1)},c^{(1)})F(h_{1,n_1}^{(2)},h_{2,n_2}^{(2)},h_{3,n_3}^{(2)},c^{(2)}).
\fe




\subsubsection*{The NS branching coefficient}

The remaining problem is to compute the branching coefficient $B_a(h_1,h_2,h_3;n_1,n_2,n_3)$. This has been worked out in \cite{1312.4520, 1406.3008} (they called it the blow up factor). To compute this factor, let us define the normalization of $v_n$ by the following. 

We first introduce a free field representation of the $\mathsf{SCA}_{1/2}$ algebra. The generators of the free field algebra $\mathcal{A}_{1/2}$ involve free boson operators $\left\{c_n | n\neq 0\right\}$, free fermion operators $\left\{\eta_r| r\in \mathbb{Z}+\frac{1}{2}\right\}$ (called $\psi_r$ in \cite{1608.02568}) and the central elements $\hat{P}$. They satisfy the algebraic relations
\ie
[c_n,c_m] =n\delta_{m+n,0}, ~\left\{\eta_r,\eta_s\right\}= \delta_{r+s,0}.
\fe
We denote by \(\mathcal{V}(P)\) the Fock representation of the free-field algebra on which the central element \(\hat P\) acts as \(P\). It is generated by acting $c_{-n}, \eta_{-r}$ on the lowest weight vector $\phi(P)$, with $\hat{P} \phi(P)= P \phi(P)$.

The free field representation of $\mathsf{SCA}_{1/2}$ is obtained by
\ie
\begin{aligned}
L_n & =\frac{1}{2} \sum_{k \neq 0, n} c_k c_{n-k}+\frac{1}{2} \sum_r r \eta_{n-r} \eta_r+\frac{i}{2}(Q n - 2 \hat{P}) c_n, \quad n \neq 0, \\
L_0 & =\sum_{k>0} c_{-k} c_k+\sum_{r>0} r \eta_{-r} \eta_r+\frac{1}{2}\left(\frac{Q^2}{4}-\hat{P}^2\right), \\
G_r & =\sum_{n \neq 0} c_n \eta_{r-n}+i(Q r - \hat{P}) \eta_r.
\end{aligned}
\fe
The Fock representation \(\mathcal{V}(P)\) is therefore also a representation of $\mathsf{SCA}_{1/2}$.
For $P\notin \frac{1}{2}(b\mathbb{Z}+b^{-1}\mathbb{Z})$, this representation is irreducible as a $\mathsf{SCA}_{1/2}$ representation and is isomorphic to the highest weight representation $\mathcal{V}_{\mathsf{SCA}_{1/2}}(h)$ with $h = \frac{1}{8}Q^2-\frac{1}{2}P^2$. 

Consider the tensor-product representation $\mathcal{M}(P)
:=
\mathcal{V}_{\mathsf{F}_{1/2}} \otimes \mathcal{V}(P)
$, where \(\mathcal{V}_{\mathsf{F}_{1/2}}\) is the Fock representation of the
auxiliary free fermion. 
Since \(\mathcal{M}(P)\)
is a representation of
\(\mathsf{F}_{1/2}\oplus\mathsf{SCA}_{1/2}\) and two copy of virasoro algebra \(\mathsf{Vir}\oplus\mathsf{Vir}\)  is a subalgebra of \(\mathsf{F}_{1/2}\oplus\mathsf{SCA}_{1/2}\),
it is a representation of \(\mathsf{Vir}\oplus\mathsf{Vir}\). 

For $P\notin \frac{1}{2}(b\mathbb{Z}+b^{-1}\mathbb{Z})$, $\mathcal{M}(P)$ is isomorphic to $\mathcal{M}_{NS}(h)$ with $h = \frac{1}{8}Q^2-\frac{1}{2}P^2$. Under this identification, \(\mathcal{M}(P)\) admits two natural
bases. One is the Fock basis inherited from $\mathcal{M}(P)$, obtained by acting the free field generators $c_{-n}, \eta_{-r} ,\psi_{-s}$ on the lowest weight vector $\mathbf{1}_{\mathsf{F}}\otimes \phi(P)$. Another basis is the PBW basis inherited from $\mathcal{M}_{NS}(h)$, obtained by acting the super-virasoro generators $L_{-n}, G_{-r} ,\psi_{-s}$ on $\mathbf{1}_{\mathsf{F}}\otimes \phi(P)$. One can thus consider the base transformation between Fock basis and PBW basis, which a prior would be complicated.

For our application, we will always endow $\mathcal{M}(P)$ with the BPZ pairing, under which $L_n^\dagger = L_{-n}, G_{r}^\dagger = G_{-r}$. Under this pairing, one can speak of the BPZ conjugation of free field generators $c_n,\eta_r$, which generically does not admit a simple formula in terms of the original free field generators. Due to this reason, it is not straightforward to compute  the norm of the state in the Fock basis. The efficient way to compute the norm is to do the base transformation to the PBW basis and then compute the norm there.
\todo[inline]{Yutai: The derivation in the NS sector yet to be reviewed. The point seems to be that the BPZ pairing is better defined as the pairing between \(\mathcal{M}(P)\) and \(\mathcal{M}(-P)\). The BPZ conjugate then maps the super-Virasoro generator in one representation to another (isomorphic) representation. It is different from the usual representation theory setup where the pairing/inner product is defined on the same space.}

The $\mathsf{Vir}\oplus\mathsf{Vir}$ primaries for $n>0$ in $\mathcal{M}(P)$ admits a simple formula in the Fock basis
\ie
w_n(P)=\chi_{-\frac{4n-1}{2}}(P)\chi_{-\frac{4n-3}{2}}(P)\cdots \chi_{-\frac{1}{2}}(P)\phi(P),\quad \chi_{r} = \psi_r - i\eta_r.
\fe
For each free boson mode $c_n$ and free fermion mode $\eta$, one can directly solve them in terms of $G_{-r}, L_{-n}$ and $\hat{P}$. We define the Virasoro primary $w_n$ in $\mathcal{M}_{\NS}(h)$ to be the state after replacing $c, \eta$ by $L, G$ modes, $\phi(P)$ by $\phi$, and $\hat{P}$ takes the positive root of $\frac{1}{4}Q^2-2h$. If $h$ is greater than $Q^2/8$, we take $P=ip$ such that $p>0$. The two states $w_n(P)$ and $w_n$ will have the same norm.

Under the BPZ inner product where $L_n^\dagger = L_{-n}, G_{r}^\dagger = G_{-r}$, we adopt the following conjugation rules for $c, \eta, \hat{P}$:
\ie
c^\dagger_n = c_{-n},\quad \eta_{r}^\dagger = \eta_{-r}, \quad \hat{P}^\dagger = -\hat{P}
\fe
and hence the BPZ conjugate of $\chi_{r}(P)$ is $-{\chi}_{-r}(-P) = -\psi_{-r}+i\eta_{-r}(-P)$. The remaining step of computing $\|w_n\|^2$ is then determining the commutator $\{\chi_r(-P), \chi_s(P)\}$. This is worked out in \cite{1312.4520} by analyzing the pole structure of this quantity in $P$, and it turns out that
\ie
\|w_n(P)\|^2 = (-1)^{2n} 2^{2n} \frac{\ell(2P,4n)}{\ell(Q+2P,4n)}
\fe
where the $\ell$ function is defined by the following. For positive integer $m$,
\ie
\ell(x, m)= \begin{cases}\prod_{\substack{r, s \geq 0, r+s<m \\ r+s \equiv 0(2)}}\left(x+r b+s b^{-1}\right), & m>0 \text { even } \\ 2^{1 / 8} \prod_{\substack{r, s \geq 0, r+s<m \\ r+s \equiv 1(2)}}\left(x+r b+s b^{-1}\right), & m>0 \text { odd }\end{cases}
\fe
For zero or negative integer $m$, 
\ie
\ell(x, 0)=1, \quad \ell(x, m)= \begin{cases}(-1)^{-m / 2} \ell(Q-x,-m), & m<0 \text { even } \\ \ell(Q-x,-m), & m<0 \text { odd }\end{cases}
\fe
This formula has been checked against direct computation of the state $w_n$ up to level $n=\frac{3}{2}$.

We normalize the Virasoro primary $v_n$ with positive $n$ to be
\ie
v_n = 2^{-2n} \ell(Q+2P, 4n)w_n,
\fe
where, again, $P$ is taken to be the positive for small $h$ and purely imaginary with positive imaginary part for large $h$. $v_0$ is defined to be $\phi$, and $v_{n}$ with negative $n$ is defined from the reflection formula
\ie
v_n(P) = v_{-n}(-P).
\fe
The norm of $v_n$ is given by
\ie
\|v_n\|^2 = (-1)^{2n}2^{-2n} \ell(2P, 4n)\ell(Q+2P, 4n).
\fe

The 3-point function $\hat{\rho}_a(v_{1,n_1},v_{2,n_2},v_{3,n_3})$ is also computed in \cite{1312.4520}. Normalizing
\ie
\hat{\rho}_0(\phi_1, \phi_2, \phi_3) = 1,\quad \hat{\rho}_1(\phi_1, G_{-1/2}\phi_2, \phi_3)=1.
\fe
In our convention, their result reads
\ie
\hat{\rho}_a&(v_{1,n_1},v_{2,n_2},v_{3,n_3}) =(-1)^{2n_3} (-1)^{n_1+n_2+n_3}i^a 2^{-\left(n_1+n_2+n_3\right)}  \\
&\quad \times \ell\left(\tfrac{Q}{2}+P_1+P_2+P_3, 2(n_1+n_2+n_3)\right)\ell\left(\tfrac{Q}{2}-P_1+P_2+P_3, 2(-n_1+n_2+n_3)\right)\\
&\quad \times \ell\left(\tfrac{Q}{2}+P_1+P_2-P_3, 2(n_1+n_2-n_3)\right) \ell\left(\tfrac{Q}{2}-P_1+P_2-P_3, 2(-n_1+n_2-n_3)\right).
\fe
For $(-1)^{2(n_1+n_2+n_3)}\neq (-1)^a$, the 3-point function is 0 by definition. Again, we should emphasize that this result only works for positive $n_i$. For negative $n_i$, one should use the reflection relation $v_n(P)=v_{-n}(-P)$ to convert it to a positive value. The values of $P_i$ is the same in the definition of $v_n$. This formula has been checked directly against low-level $(n=0,1,\frac{1}{2})$ 3-point functions.

Putting everything together, we arrive at the final form for the NS branching coefficient:
\ie
B_a&(h_1, h_2, h_3;n_1, n_2, n_3) = \frac{(-1)^{2n_3}(-i)^a}{\prod_{i=1}^3\sqrt{\ell(2P_i, 4n_i)\ell(Q+2P_i, 4n_i)}}\\
&\quad \times \ell\left(\tfrac{Q}{2}+P_1+P_2+P_3, 2(n_1+n_2+n_3)\right)\ell\left(\tfrac{Q}{2}-P_1+P_2+P_3, 2(-n_1+n_2+n_3)\right)\\
&\quad \times \ell\left(\tfrac{Q}{2}+P_1+P_2-P_3, 2(n_1+n_2-n_3)\right) \ell\left(\tfrac{Q}{2}-P_1+P_2-P_3, 2(-n_1+n_2-n_3)\right).
\fe


\subsubsection*{Checks against $c$-recursion}

As a check of the formalism developed in this section, we compute the all-NS block $\mathbf{F}_a(h_1, h_2, h_3, c)$ as a $q$-expansion independently using the $c$-recursion relation and the double Virasoro decomposition. For all the vacuum blocks appearing in $\mathbf{F},\mathbf{F}_{\mathsf{F}}$ and $F$, we use the Schottky frame so that the Weyl anomaly factor on both sides agrees. We have tested numerically that, up to total level 4, the all-NS theta channel block $\mathbf{F}_a$ computed from the two methods agrees to machine precision for all 8 independent choices of $\{\eta_i\}$ and $a$. 

\section{Ramond blocks from double Virasoro}

We now generalize the discussion above to Ramond blocks. Denoting the fermion ground states by $u^\mathsf{a}$ with $\mathsf{a}=0,1$, normalizing
\ie
\langle u^0|u^0\rangle = 1,\quad \langle u^1|u^1 \rangle=-1,\quad \psi_0 u^\mathsf{a}=\frac{1}{\sqrt{2}}u^{1-\mathsf{a}},
\fe
the ramond $\mathsf{SCA}_0\otimes \mathsf{F}_0$ block is defined as
\ie 
\widehat{\mathbb{F}}_f^{(\eta, \eta')}&(h_1, \beta_2, \beta_3, c) = \sum_{
\substack{|A|=|A'|, |B|=|B'|, |C|=|C'|\\
(-1)^{A+B+C+|\alpha|+|\gamma|}=(-1)^f\\
(-1)^{A'+B'+C'+|\alpha'|+|\gamma'|}=(-1)^f\\
(-1)^{\mathsf{A+B+C+b+c}}=1}
} q_1^{|A|} \eta_1^A q_2^{|B|} \eta_2^{B+|\alpha|} q_3^{|C|} \eta_3^{C+|\gamma|} \mathbf{G}^{AA'}_{h_1} \mathbb{G}^{B\alpha, B'\alpha'}_{\beta_2} \mathbb{G}^{C\gamma, C'\gamma'}_{\beta_3}\\
&\quad\quad \times (-1)^{\text{sgn}(A)\text{sgn}(B)+\text{sgn}(A)\text{sgn}(C)+\text{sgn}(B)\text{sgn}(C)} q_1^{|\mathsf{A}|} \eta_1^{\mathsf{A}} q_2^{|\mathsf{B}|} \eta_2^{\mathsf{B}+\mathsf{b}} q_3^{|\mathsf{C}|} \eta_3^{\mathsf{C}+\mathsf{c}}\\
&\quad\quad \times \hat{\rho}_f^{(\eta)}(\Psi_{-\mathsf{A}}\mathbf{L}_{-A}\phi_1, \boldsymbol\Psi_{-\mathsf{B}}u^{\mathsf{b}}\otimes \mathbb{L}_{-B}w^\alpha_2, \boldsymbol\Psi_{-\mathsf{C}}u^{\mathsf{c}}\otimes \mathbb{L}_{-C}w^\gamma_3)\\
&\quad\quad \times \hat{\rho}_f^{(\eta')}(\Psi_{-\mathsf{A}}\mathbf{L}_{-A'}\phi_1, \boldsymbol\Psi_{-\mathsf{B}}u^{\mathsf{b}}\otimes \mathbb{L}_{-B'}w^{\alpha'}_2, \boldsymbol\Psi_{-\mathsf{C}}u^{\mathsf{c}} \otimes\mathbb{L}_{-C'}w^{\gamma'}_3)
\fe
where $\boldsymbol{\Psi}_{-\mathsf{B}}$ denotes a chain of R-sector fermion lowering operators (there are no \texttt{mathbb} version of $\Psi$, sadly). Note that $\boldsymbol\Psi_{-\mathsf{A}}$ does not contain the fermion zero mode $\psi_0$. We have used $\langle u^\mathsf{a}|\boldsymbol\Psi_{-\mathsf{A}}^\dagger \boldsymbol\Psi_{-\mathsf{B}}|u^{\mathsf{b}}\rangle=(-1)^{\mathsf{A}+\mathsf{a}}\delta_{\mathsf{AB}}\delta^{\mathsf{a}\mathsf{b}}$, as well as $\langle u_1\otimes x_1| u_2\otimes x_2 \rangle=\langle u_1|u_2\rangle \langle x_1|x_2\rangle$ without the additional operator exchanging sign. The shorthand $\text{sgn}$ is defined such that ${\text{sgn}(A)}=A+\mathsf{A} \mod 2$ for NS sector and ${\text{sgn}(B)}=B+\mathsf{B}+\mathsf{b}+|\alpha| \mod 2$ for R sector. We have also used the fact that the fermion correlator is 0 unless $(-1)^{\mathsf{A+B+C+b+c}}=1$.

The 3-point function can be factorized as (we treat this formula as the definition of $\hat{\rho}^{(\eta)}_f$)
\ie
\hat{\rho}_f^{(\eta)}&(\Psi_{-\mathsf{A}}\mathbf{L}_{-A}\phi_1, \boldsymbol\Psi_{-\mathsf{B}}u^{\mathsf{b}}\otimes \mathbb{L}_{-B}w^\alpha_2, \boldsymbol\Psi_{-\mathsf{C}}u^{\mathsf{c}}\otimes \mathbb{L}_{-C}w^\gamma_3)\\
&=(-1)^{A\mathsf{A}+(B+|\alpha|)(\mathsf{C+c})+f(\mathsf{C+c})}\rho(\Psi_{-\mathsf{A}}\mathbf{1}_{\mathsf{F}}, \boldsymbol\Psi_{-\mathsf{B}}u^{\mathsf{b}},\boldsymbol\Psi_{-\mathsf{C}}u^{\mathsf{c}}) {\rho}_f^{(\eta)}(\mathbf{L}_{-A}\phi_1, \mathbb{L}_{-B}w^\alpha_2, \mathbb{L}_{-C}w^\gamma_3) 
\fe
Therefore, if we define the fermion Ramond block as
\ie
\mathbb{F}_{\mathsf{F}} = \sum_{\mathsf{A+B+C+b+c}\equiv 0 \mod 2} &q_1^{|\mathsf{A}|} \eta_1^{\mathsf{A}} q_2^{|\mathsf{B}|} \eta_2^{\mathsf{B}+\mathsf{b}} q_3^{|\mathsf{C}|} \eta_3^{\mathsf{C}+\mathsf{c}} (-1)^{\mathsf{A(B+b)+A(C+c)+(B+b)(C+c)}}\\
& \times [\rho(\Psi_{-\mathsf{A}}\mathbf{1}_{\mathsf{F}}, \boldsymbol\Psi_{-\mathsf{B}}u^{\mathsf{b}},\boldsymbol\Psi_{-\mathsf{C}}u^{\mathsf{c}})]^2,
\fe
one then have the decomposition
\[\widehat{\mathbb{F}}_f^{(\eta, \eta')}(h_1, \beta_2, \beta_3, c) = \mathbb{F}_{\mathsf{F}}\star_{\R}\mathbb{F}_f^{(\eta, \eta')}(h_1, \beta_2, \beta_3, c)\]
where the R-sector convolution acts on monomials of $q_i$ and $\eta_i$ as
\ie
q_1^{r_1}q_2^{r_2}q_3^{r_3}&\eta_1^{\epsilon_1}\eta_2^{\epsilon_2}\eta_3^{\epsilon_3} \star_{\R} q_1^{s_1}q_2^{s_2}q_3^{s_3}\eta_1^{\delta_1}\eta_2^{\delta_2}\eta_3^{\delta_3}\\
&=(-1)^{\sum_{i<j} (\epsilon_i \delta_j + \delta_i \epsilon_j)}q_1^{r_1+s_1}q_2^{r_2+s_2}q_3^{r_3+s_3}\eta_1^{\epsilon_1+\delta_1}\eta_2^{\epsilon_2+\delta_2}\eta_3^{\epsilon_3+\delta_3}.
\fe

Using the decomposition \eqref{Rbranching}, one can see that every state in $\mathcal{M}(\beta)$ can be written as $L_{-A}^{(1)}L_{-B}^{(2)}v^\alpha_n$, with $\alpha = \pm$ and $n\in \frac{1}{2}\mathbb{Z}+\frac{1}{4}$. The state $v^\alpha_n$ has conformal weight $h_n^{(1)}+h_{n}^{(2)}-h=2n^2 -\frac{1}{16}$ and fermion parity $(-1)^\alpha$. Therefore,
\ie
\widehat{\mathbb{F}}^{(\eta, \eta')}_f&(h_1, \beta_2, \beta_3, c) = \sum_{\substack{
  |A|=|A'|, |B|=|B'|, |C|=|C'|\\
  |D|=|D'|, |E|=|E'|, |F|=|F'|\\
  n_1\in \frac{1}{2}\mathbb{Z},~ n_2, n_3\in \frac{1}{2}\mathbb{Z}+\frac{1}{4}\\
  (-1)^{2n_1 + \alpha_2 + \alpha_3}=(-1)^f
  }} q_1^{2n_1^2+|A|+|B|}q_2^{2n_2^2-\frac{1}{8}+|C|+|D|}q_3^{2n_3^2-\frac{1}{8}+|E|+|F|}\eta_1^{2n_1}\eta_2^{\alpha_2}\eta_3^{\alpha_3}\\
&\times (-1)^{2n_1\alpha_2+2n_1\alpha_3+\alpha_2\alpha_3} \frac{(G_{h_{1,n_1}^{(1)}})^{AA'}(G_{h_{1,n_1}^{(2)}})^{BB'}(G_{h_{2,n_2}^{(1)}})^{CC'}(G_{h_{2,n_2}^{(2)}})^{DD'}(G_{h_{3,n_3}^{(1)}})^{EE'}(G_{h_{3,n_3}^{(2)}})^{FF'}}{\|v_{1,n_1}\|^2\|v_{2,n_2}^{\alpha_2}\|^2\|v_{3,n_3}^{\alpha_3}\|^2}\\
&\times \hat{\rho}_f^{(\eta)}(L_{-A}^{(1)}L_{-B}^{(2)}v_{1,n_1},L_{-C}^{(1)}L_{-D}^{(2)}v_{2,n_2}^{\alpha_2},L_{-E}^{(1)}L_{-F}^{(2)}v_{3,n_3}^{\alpha_3})\hat{\rho}_f^{(\eta')}(L_{-A'}^{(1)}L_{-B'}^{(2)}v_{1,n_1},L_{-C'}^{(1)}L_{-D'}^{(2)}v_{2,n_2}^{\alpha_2},L_{-E'}^{(1)}L_{-F'}^{(2)}v_{3,n_3}^{\alpha_3}),
\fe
where we used the fact that $u^{\mathsf{a}}$ has conformal weight $\frac{1}{16}$ and the fact that $\langle v_n^0|v_n^1\rangle = 0$. The factorization relation of the 3-point function reads
\ie
\hat{\rho}_f^{(\eta)}&(L_{-A}^{(1)}L_{-B}^{(2)}v_{1,n_1},L_{-C}^{(1)}L_{-D}^{(2)}v_{2,n_2}^{\alpha_2},L_{-E}^{(1)}L_{-F}^{(2)}v_{3,n_3}^{\alpha_3})\\
&\quad =\hat{\rho}_f^{(\eta)}(v_{1,n_1}, v_{2,n_2}^{\alpha_2}, v_{3,n_3}^{\alpha_3})\rho(L_{-A}\nu_{h_{1,n_1}^{(1)}}, L_{-C}\nu_{h_{2,n_2}^{(1)}}, L_{-E}\nu_{h_{3,n_3}^{(1)}})\rho(L_{-B}\nu_{h_{1,n_1}^{(2)}}, L_{-D}\nu_{h_{2,n_2}^{(2)}}, L_{-F}\nu_{h_{3,n_3}^{(2)}}),
\fe
and therefore,
\ie
\widehat{\mathbb{F}}_f^{(\eta,\eta')}(h_1, \beta_2, \beta_3, c) &= \sum_{\substack{n_1 \in \frac{1}{2} \mathbb{Z}, n_2, n_3\in \frac{1}{2}\mathbb{Z}+\frac{1}{4}\\(-1)^{2n_1+\alpha_2+\alpha_3}=(-1)^f}} q_1^{2n_1^2}q_2^{2n_2^2-\frac{1}{8}}q_3^{2n_3^2-\frac{1}{8}}\eta_1^{2n_1}\eta_2^{\alpha_2}\eta_3^{\alpha_3}\\
&\quad \times (-1)^{2n_1\alpha_2+2n_1\alpha_3+\alpha_2\alpha_3} \mathbb{B}^{(\eta)}_f(h_1, \beta_2, \beta_3;n_1, n_2, n_3) \mathbb{B}^{(\eta')}_f(h_1, \beta_2, \beta_3;n_1, n_2, n_3)\\
&\quad \times F(h_{1,n_1}^{(1)},h_{2,n_2}^{(1)},h_{3,n_3}^{(1)},c^{(1)})F(h_{1,n_1}^{(2)},h_{2,n_2}^{(2)},h_{3,n_3}^{(2)},c^{(2)}),
\fe
where the R-sector branching coefficients are defined as
\ie
\mathbb{B}^{(\eta)}_f(h_1, \beta_2, \beta_3;n_1, n_2, n_3) = \frac{\hat{\rho}_f^{(\eta)}(v_{1,n_1}, v_{2,n_2}^{\alpha_2}, v_{3,n_3}^{\alpha_3})}{\|v_{1,n_1}\|\|v_{2,n_2}^{\alpha_2}\|\|v_{3,n_3}^{\alpha_3}\|}.
\fe

\newpage
\begin{appendix}
\section{Fusion polynomials and factorization formulae}

We collect the result for all useful fusion polynomials and factorization formulae in this appendix. They are taken directly from the menu of a fusion restaurant in Cambridge, MA. All the formulas here are tested against direct low-level computations and checked to be consistent with the Mobius tranformation that permutes the three slots.

\subsection{Fusion polynomials on a (NS, NS, NS) sphere}

The All-NS fusion polynomials are defined as
\ie
\mathbf{P}_{rs}^{0}(h_i,h_j;c)
&={}&\prod_{\substack{
p=1-r,1-r+2,\ldots,r-1\\
q=1-s,1-s+2,\ldots,s-1\\
p+q-(r+s)\in4\mathbb Z+2}}
\frac{\lambda_i-\lambda_j+pb+qb^{-1}}{2\sqrt2}
\frac{\lambda_i+\lambda_j+pb+qb^{-1}}{2\sqrt2},\\
\mathbf{P}_{rs}^{1}(h_i,h_j;c)
&={}&\prod_{\substack{
p=1-r,1-r+2,\ldots,r-1\\
q=1-s,1-s+2,\ldots,s-1\\
p+q-(r+s)\in4\mathbb Z}}
\frac{\lambda_i-\lambda_j+pb+qb^{-1}}{2\sqrt2}
\frac{\lambda_i+\lambda_j+pb+qb^{-1}}{2\sqrt2}.
\label{eq:ns-fusion-polynomials}
\fe
They satisfy the reflection relation
\ie
 \mathbf{P}_{rs}^{0}(h_i,h_j;c_{rs})
 &= (-1)^{\lceil rs/2 \rceil}
   \mathbf{P}_{rs}^{0}(h_j,h_i;c_{rs}),\\
 \mathbf{P}_{rs}^{1}(h_i,h_j;c_{rs})
 &=(-1)^{\lfloor rs/2\rfloor}
   \mathbf{P}_{rs}^{1}(h_j,h_i;c_{rs}).
 \label{eq:appendix-ns-polynomial-reflection}
\fe
For a null state in the first slot, the factorization relation is given by
\ie
&\rho_a(\mathbf{L}_{-A}\chi_{rs} , \mathbf{L}_{-B}\phi_2,  \mathbf{L}_{-C}\phi_3) \\& \quad= \begin{cases} \rho_{a}(\mathbf{L}_{-A}\phi_{h_{rs}+\frac{rs}{2}} , \mathbf{L}_{-B}\phi_2,  \mathbf{L}_{-C}\phi_3) \mathbf{P}_{rs}^a(h_3, h_2),\quad rs\in 2\mathbb{Z}, \\ 
(-1)^{p_1+A}\rho_{1-a}(\mathbf{L}_{-A}\phi_{h_{rs}+\frac{rs}{2}} , \mathbf{L}_{-B}\phi_2,  \mathbf{L}_{-C}\phi_3) \mathbf{P}_{rs}^{a}(h_3, h_2),\quad rs\in 2\mathbb{Z}+1,
\end{cases}
\label{eq:appendix-ns-first-factorization}
\fe
For a null state in the second slot, 
\ie
&\rho_a\!\left(
\mathbf L_{-A}\phi_1,
\mathbf L_{-C}\chi_{rs},
\mathbf L_{-E}\phi_3
\right)
\\
&\quad=
\begin{cases}
\mathbf P_{rs}^{a}(h_3,h_1;c_{rs})\rho_a\!\left(
\mathbf L_{-A}\phi_1,
\mathbf L_{-C}\phi_{h_{rs}+\frac{rs}{2}},
\mathbf L_{-E}\phi_3
\right),
& rs\in2\mathbb Z,
\\[0.4em]
\mathbf P_{rs}^{a}(h_3,h_1;c_{rs})\rho_{1-a}\!\left(
\mathbf L_{-A}\phi_1,
\mathbf L_{-C}\phi_{h_{rs}+\frac{rs}{2}},
\mathbf L_{-E}\phi_3
\right),
& rs\in2\mathbb Z+1.
\end{cases}
\label{eq:appendix-ns-second-factorization}
\fe
and for the third slot,
\ie
&\rho_a\!\left(
\mathbf L_{-A}\phi_1,
\mathbf L_{-C}\phi_2,
\mathbf L_{-E}\chi_{rs}
\right)
\\
&\quad=
\begin{cases}
\mathbf P_{rs}^{a}(h_1,h_2;c_{rs})\rho_a\!\left(
\mathbf L_{-A}\phi_1,
\mathbf L_{-C}\phi_2,
\mathbf L_{-E}\phi_{h_{rs}+\frac{rs}{2}}
\right),
& rs\in2\mathbb Z,
\\[0.4em]
(-1)^{1+p_2}
\mathbf P_{rs}^{a}(h_1,h_2;c_{rs})\rho_{1-a}\!\left(
\mathbf L_{-A}\phi_1,
\mathbf L_{-C}\phi_2,
\mathbf L_{-E}\phi_{h_{rs}+\frac{rs}{2}}
\right),
& rs\in2\mathbb Z+1.
\end{cases}
\label{eq:appendix-ns-third-factorization}
\fe
For the case where two of the three states are null states, one have
\ie
\rho_f(\mathbf{L}_{-A}\chi_{rs},
        \mathbf{L}_{-C}\chi_{rs},
        \mathbf{L}_{-E}\phi_3)&=(-1)^{rs(p_1+A)}
 \mathbf{P}_{rs}^{f}(h_3,h_{rs};c_{rs})
 \mathbf{P}_{rs}^{f+rs}
 \left(h_3,h_{rs}+\frac{rs}{2};c_{rs}\right)\\
&\qquad\quad\times
 \rho_f(\mathbf{L}_{-A}\phi_{h_{rs}+rs/2},
        \mathbf{L}_{-C}\phi_{h_{rs}+rs/2},
        \mathbf{L}_{-E}\phi_3),
\\
\rho_f(\mathbf{L}_{-A}\chi_{rs},
        \mathbf{L}_{-C}\phi_2,
        \mathbf{L}_{-E}\chi_{rs})&=(-1)^{rs(1+p_1+p_2+A)}
 \mathbf{P}_{rs}^{f}(h_{rs},h_2;c_{rs})
 \mathbf{P}_{rs}^{f+rs}
 \left(h_{rs}+\frac{rs}{2},h_2;c_{rs}\right)\\
&\qquad\quad\times
 \rho_f(\mathbf{L}_{-A}\phi_{h_{rs}+rs/2},
        \mathbf{L}_{-C}\phi_2,
        \mathbf{L}_{-E}\phi_{h_{rs}+rs/2}),
\\
\rho_f(\mathbf{L}_{-A}\phi_1,
        \mathbf{L}_{-C}\chi_{rs},
        \mathbf{L}_{-E}\chi_{rs})&=(-1)^{rsp_2}
 \mathbf{P}_{rs}^{f}(h_{rs},h_1;c_{rs})
 \mathbf{P}_{rs}^{f+rs}
 \left(h_1,h_{rs}+\frac{rs}{2};c_{rs}\right)\\
&\qquad\quad\times
 \rho_f(\mathbf{L}_{-A}\phi_1,
        \mathbf{L}_{-C}\phi_{h_{rs}+rs/2},
        \mathbf{L}_{-E}\phi_{h_{rs}+rs/2}).
\fe

\subsection{NS null states on a (NS, R, R) sphere}
The two fusion polynomial for a NS null state on a (NS, R, R) sphere is defined as
\ie
\mathbb{P}_{rs}^{\NS,\eta}(\beta_2,\beta_3;c)
={}&
\prod_{\substack{0\le k\le r-1,\ 0\le\ell\le s-1\\k+\ell\in2\mathbb Z}}
\frac{2\sqrt2(\beta_3-\eta\beta_2)
 -(1-r+2k)b-(1-s+2\ell)b^{-1}}{2\sqrt2}\\
&\times
\prod_{\substack{0\le k\le r-1,\ 0\le\ell\le s-1\\k+\ell\in2\mathbb Z+1}}
\frac{2\sqrt2(\beta_3+\eta\beta_2)
 -(1-r+2k)b-(1-s+2\ell)b^{-1}}{2\sqrt2}.
\label{eq:ns-rr-fusion-polynomial}
\fe
For the NS null state in the first slot, the factorization relation reads
\ie
&\rho_f^{(\eta)}
 (\mathbf{L}_{-A}\chi_{rs},
  \mathbb{L}_{-C}w_{\beta_2}^{\alpha},
  \mathbb{L}_{-E}w_{\beta_3}^{\gamma})\\
&\quad=\mathbb{P}_{rs}^{\NS,\eta}(\beta_2,\beta_3;c_{rs})
\begin{cases}
\rho_f^{(\eta)}
 (\mathbf{L}_{-A}\phi_{h_{rs}+rs/2},
  \mathbb{L}_{-C}w_{\beta_2}^{\alpha},
  \mathbb{L}_{-E}w_{\beta_3}^{\gamma}),&rs\in2\mathbb Z,\\[0.4em]
e^{i(-1)^f\pi/4}\rho_{1-f}^{(-\eta)}
 (\mathbf{L}_{-A}\phi_{h_{rs}+rs/2},
  \mathbb{L}_{-C}w_{\beta_2}^{\alpha},
  \mathbb{L}_{-E}w_{\beta_3}^{\gamma}),&rs\in2\mathbb Z+1.
\end{cases}
\fe
For the third slot,
\ie
&\rho_f^{(\eta)}
 (\mathbb{L}_{-A}w_{\beta_1}^{\alpha},
  \mathbb{L}_{-C}w_{\beta_2}^{\gamma},
  \mathbf{L}_{-E}\chi_{rs})
\\
&\quad=
\mathbb{P}_{rs}^{\NS,(-1)^f\eta}(\beta_2,\beta_1;c_{rs})
\begin{cases}
\rho_f^{(\eta)}
 (\mathbb{L}_{-A}w_{\beta_1}^{\alpha},
  \mathbb{L}_{-C}w_{\beta_2}^{\gamma},
  \mathbf{L}_{-E}\phi_{h_{rs}+rs/2}),&rs\in2\mathbb Z,\\[0.4em]
\eta e^{i(-1)^f\pi/4}\rho_{1-f}^{(\eta)}
 (\mathbb{L}_{-A}w_{\beta_1}^{\alpha},
  \mathbb{L}_{-C}w_{\beta_2}^{\gamma},
  \mathbf{L}_{-E}\phi_{h_{rs}+rs/2}),&rs\in2\mathbb Z+1.
\end{cases}
\fe
Finally, for the second slot with even $rs$, one have
\ie
&\rho_f^{(\eta)}
 (\mathbb{L}_{-A}w_{\beta_1}^{\alpha},
  \mathbf{L}_{-C}\chi_{rs},
  \mathbb{L}_{-E}w_{\beta_3}^{\gamma})
\\
&\quad=\frac12\sum_{\epsilon=\pm1}
 \mathbb{P}_{rs}^{\NS,\epsilon}(\beta_3,\beta_1;c_{rs})
 \Big[
 \rho_f^{(\eta)}
 (\mathbb{L}_{-A}w_{\beta_1}^{\alpha},
  \mathbf{L}_{-C}\phi_{h_{rs}+rs/2},
  \mathbb{L}_{-E}w_{\beta_3}^{\gamma})
\\[-0.2em]
&\hspace{55mm}
 {}-i\eta\epsilon\,
 \rho_f^{(-\eta)}
 (\mathbb{L}_{-A}w_{\beta_1}^{\alpha},
  \mathbf{L}_{-C}\phi_{h_{rs}+rs/2},
  \mathbb{L}_{-E}w_{\beta_3}^{\gamma})
 \Big].
\label{eq:appendix-rnr-ns-even-factorization}
\fe
and for odd $rs$, 
\ie
&\rho_f^{(\eta)}
 (\mathbb{L}_{-A}w_{\beta_1}^{\alpha},
  \mathbf{L}_{-C}\chi_{rs},
  \mathbb{L}_{-E}w_{\beta_3}^{\gamma})
\\
&\quad=\frac{(-1)^{(rs-1)/2}i^f e^{i\pi/4}}2
 \sum_{\epsilon=\pm1}
 \mathbb{P}_{rs}^{\NS,\epsilon}(\beta_3,\beta_1;c_{rs})
 \Big[
 \eta\,\rho_{1-f}^{(\eta)}
 (\mathbb{L}_{-A}w_{\beta_1}^{\alpha},
  \mathbf{L}_{-C}\phi_{h_{rs}+rs/2},
  \mathbb{L}_{-E}w_{\beta_3}^{\gamma})
\\[-0.2em]
&\hspace{70mm}
 {}+i\epsilon\,
 \rho_{1-f}^{(-\eta)}
 (\mathbb{L}_{-A}w_{\beta_1}^{\alpha},
  \mathbf{L}_{-C}\phi_{h_{rs}+rs/2},
  \mathbb{L}_{-E}w_{\beta_3}^{\gamma})
 \Big].
 \fe

\subsection{Ramond null state on a (NS, R, R) sphere}

The two fusion polynomials associated with a Ramond null state are defined as 
\ie
\mathbb{P}_{rs}^{\R,\eta}(h_i,\beta_j;c)
={}&
\prod_{\substack{0\le k\le r-1,\ 0\le\ell\le s-1\\k+\ell\in2\mathbb Z}}
\frac{\lambda_i-2\sqrt2\,\eta\beta_j
 -(1-r+2k)b-(1-s+2\ell)b^{-1}}{2\sqrt2}\\
&\times
\prod_{\substack{0\le k\le r-1,\ 0\le\ell\le s-1\\k+\ell\in2\mathbb Z+1}}
\frac{\lambda_i+2\sqrt2\,\eta\beta_j
 -(1-r+2k)b-(1-s+2\ell)b^{-1}}{2\sqrt2}.
\label{eq:r-fusion-polynomial}
\fe
The factorization relations are
\ie
&\rho_f^{(\eta)}
 (\mathbf{L}_{-A}\phi_i,
  \mathbb{L}_{-C}\chi_{rs}^{\alpha},
  \mathbb{L}_{-E}w_{\beta_j}^{\gamma})=
 \mathbb{P}_{rs}^{\R,\eta}(h_i,\beta_j;c_{rs}^{\R})
 \rho_f^{(\eta)}
 (\mathbf{L}_{-A}\phi_i,
  \mathbb{L}_{-C}w_{\beta'_{rs}}^{\alpha},
  \mathbb{L}_{-E}w_{\beta_j}^{\gamma}),
\\[0.6em]
&\rho_f^{(\eta)}
 (\mathbf{L}_{-A}\phi_i,
  \mathbb{L}_{-C}w_{\beta_j}^{\gamma},
  \mathbb{L}_{-E}\chi_{rs}^{\alpha})=
 \mathbb{P}_{rs}^{\R,\eta}(h_i,\beta_j;c_{rs}^{\R})
 \rho_f^{(\eta)}
 (\mathbf{L}_{-A}\phi_i,
  \mathbb{L}_{-C}w_{\beta_j}^{\gamma},
  \mathbb{L}_{-E}w_{\beta'_{rs}}^{\alpha}),
\\[0.6em]
&\rho_f^{(\eta)}
 (\mathbb{L}_{-A}\chi_{rs}^{\alpha},
  \mathbb{L}_{-C}w_{\beta_j}^{\gamma},
  \mathbf{L}_{-E}\phi_i)=
 \mathbb{P}_{rs}^{\R,(-1)^f\eta}(h_i,\beta_j;c_{rs}^{\R})
 \rho_f^{(\eta)}
 (\mathbb{L}_{-A}w_{\beta'_{rs}}^{\alpha},
  \mathbb{L}_{-C}w_{\beta_j}^{\gamma},
  \mathbf{L}_{-E}\phi_i),
\\[0.6em]
&\rho_f^{(\eta)}
 (\mathbb{L}_{-A}w_{\beta_j}^{\gamma},
  \mathbb{L}_{-C}\chi_{rs}^{\alpha},
  \mathbf{L}_{-E}\phi_i)=
 \mathbb{P}_{rs}^{\R,(-1)^f\eta}(h_i,\beta_j;c_{rs}^{\R})
 \rho_f^{(\eta)}
 (\mathbb{L}_{-A}w_{\beta_j}^{\gamma},
  \mathbb{L}_{-C}w_{\beta'_{rs}}^{\alpha},
  \mathbf{L}_{-E}\phi_i),
\label{eq:appendix-nrr-rrn-r-factorizations}
\fe
and for the (R, NS, R) order,
\ie
&\rho_f^{(\eta)}
 (\mathbb{L}_{-A}\chi_{rs}^{\alpha},
  \mathbf{L}_{-C}\phi_i,
  \mathbb{L}_{-E}w_{\beta_j}^{\gamma})
\\[-0.2em]
&\quad=\frac{(-1)^{rs/2}}2
 \sum_{\epsilon=\pm1}
 \mathbb{P}_{rs}^{\R,\epsilon}(h_i,\beta_j;c_{rs}^{\R})
 \left[
 \begin{aligned}
 &
 \rho_f^{(\eta)}
 (\mathbb{L}_{-A}w_{\beta'_{rs}}^{\alpha},
  \mathbf{L}_{-C}\phi_i,
  \mathbb{L}_{-E}w_{\beta_j}^{\gamma})
\\[-0.2em]
&
 {}-i\eta\epsilon\,
 \rho_f^{(-\eta)}
 (\mathbb{L}_{-A}w_{\beta'_{rs}}^{\alpha},
  \mathbf{L}_{-C}\phi_i,
  \mathbb{L}_{-E}w_{\beta_j}^{\gamma})
 \end{aligned}
 \right],
\\[0.6em]
&\rho_f^{(\eta)}
 (\mathbb{L}_{-A}w_{\beta_j}^{\gamma},
  \mathbf{L}_{-C}\phi_i,
  \mathbb{L}_{-E}\chi_{rs}^{\alpha})
\\[-0.2em]
&\quad=\frac{(-1)^{rs/2}}2
 \sum_{\epsilon=\pm1}
 \mathbb{P}_{rs}^{\R,\epsilon}(h_i,\beta_j;c_{rs}^{\R})
 \left[
 \begin{aligned}
 &
 \rho_f^{(\eta)}
 (\mathbb{L}_{-A}w_{\beta_j}^{\gamma},
  \mathbf{L}_{-C}\phi_i,
  \mathbb{L}_{-E}w_{\beta'_{rs}}^{\alpha})
\\[-0.2em]
&
 {}-i\eta\epsilon\,
 \rho_f^{(-\eta)}
 (\mathbb{L}_{-A}w_{\beta_j}^{\gamma},
  \mathbf{L}_{-C}\phi_i,
  \mathbb{L}_{-E}w_{\beta'_{rs}}^{\alpha})
 \end{aligned}
 \right].
\label{eq:appendix-rnr-r-factorizations}
\fe

For the case where the two Ramond states are simutaneously null state,
\ie
&\rho_f^{(\eta)}
 (\mathbf{L}_{-A}\phi_i,
  \mathbb{L}_{-C}\chi_{rs}^{\alpha},
  \mathbb{L}_{-E}\chi_{rs}^{\gamma})
\\[-0.2em]
&\quad=
 \mathbb{P}_{rs}^{\R,\eta}(h_i,\beta_{rs};c_{rs}^{\R})
 \mathbb{P}_{rs}^{\R,\eta}(h_i,\beta_{rs}^{\prime};c_{rs}^{\R})
 \rho_f^{(\eta)}
 (\mathbf{L}_{-A}\phi_i,
  \mathbb{L}_{-C}w_{\beta_{rs}^{\prime}}^{\alpha},
  \mathbb{L}_{-E}w_{\beta_{rs}^{\prime}}^{\gamma}),
\\[0.6em]
&\rho_f^{(\eta)}
 (\mathbb{L}_{-A}\chi_{rs}^{\alpha},
  \mathbb{L}_{-C}\chi_{rs}^{\gamma},
  \mathbf{L}_{-E}\phi_i)
\\[-0.2em]
&\quad=
 \mathbb{P}_{rs}^{\R,(-1)^f\eta}(h_i,\beta_{rs};c_{rs}^{\R})
 \mathbb{P}_{rs}^{\R,(-1)^f\eta}(h_i,\beta_{rs}^{\prime};c_{rs}^{\R})
 \rho_f^{(\eta)}
 (\mathbb{L}_{-A}w_{\beta_{rs}^{\prime}}^{\alpha},
  \mathbb{L}_{-C}w_{\beta_{rs}^{\prime}}^{\gamma},
  \mathbf{L}_{-E}\phi_i),
\\[0.6em]
&\rho_f^{(\eta)}
 (\mathbb{L}_{-A}\chi_{rs}^{\alpha},
  \mathbf{L}_{-C}\phi_i,
  \mathbb{L}_{-E}\chi_{rs}^{\gamma})
\\[-0.2em]
&\quad=\frac12\sum_{\epsilon=\pm1}
 \mathbb{P}_{rs}^{\R,\epsilon}(h_i,\beta_{rs};c_{rs}^{\R})
 \mathbb{P}_{rs}^{\R,\epsilon}(h_i,\beta_{rs}^{\prime};c_{rs}^{\R})
 \left[
 \begin{aligned}
 &
 \rho_f^{(\eta)}
 (\mathbb{L}_{-A}w_{\beta_{rs}^{\prime}}^{\alpha},
  \mathbf{L}_{-C}\phi_i,
  \mathbb{L}_{-E}w_{\beta_{rs}^{\prime}}^{\gamma})
\\[-0.2em]
&
 \quad-i\eta\epsilon\,
 \rho_f^{(-\eta)}
 (\mathbb{L}_{-A}w_{\beta_{rs}^{\prime}}^{\alpha},
  \mathbf{L}_{-C}\phi_i,
  \mathbb{L}_{-E}w_{\beta_{rs}^{\prime}}^{\gamma})
 \end{aligned}
 \right].
\fe

\end{appendix}

\end{document}
\bye
